\documentclass[11pt]{article}

\usepackage[utf8]{inputenc}
\usepackage[T1]{fontenc}
\usepackage{natbib}
\usepackage{geometry}
\geometry{margin=2.5cm}
\usepackage{hyperref}
\usepackage{url}
\usepackage{booktabs}
\usepackage{amsmath}
\usepackage{amsfonts}
\usepackage{amssymb}
\usepackage{nicefrac}
\usepackage{microtype}
\usepackage{xcolor}
\usepackage{graphicx}
\usepackage{float}
\usepackage[table]{xcolor}
\usepackage{subcaption}
\usepackage{cleveref}
\crefname{equation}{Eq.\@}{Eqs.\@}
\Crefname{equation}{Eq.\@}{Eqs.\@}
\crefname{figure}{Fig.\@}{Figs.\@}
\Crefname{figure}{Fig.\@}{Figs.\@}
\crefname{table}{Tab.\@}{Tabs.\@}
\Crefname{table}{Tab.\@}{Tabs.\@}
\crefname{section}{Sect.\@}{Sects.\@}
\Crefname{section}{Sect.\@}{Sects.\@}

\definecolor{cgreen}{HTML}{2ea043}
\definecolor{corange}{HTML}{d29922}
\definecolor{cred}{HTML}{cf222e}
\newcommand{\good}[1]{\textcolor{green!50!black}{#1}}
\newcommand{\bad}[1]{\textcolor{orange}{#1}}

\def\eg{{\it e.g.}}
\def\ie{{\it i.e.}}
\def\R{\mathbb{R}}
\def\relu{\mathrm{ReLU}}
\def\softplus{\mathrm{softplus}}
\def\petrucco{Petrucco et al.\ 2023}
\def\rpi{r1\ensuremath{\pi}}
% Functional-class macros (Results section)
\newcommand{\mi}{\mathcal{I}}
\newcommand{\omag}{\mathcal{W}}
\newcommand{\smod}{\mathcal{S}}
\newcommand{\specf}{\mathcal{T}}
\newcommand{\cls}[1]{\mathsf{#1}}

% Figure directory. Change ONLY this line when moving between editors:
% \newcommand{\figdir}{figures/zebrafish}
%   local repo (docs/ build):              \newcommand{\figdir}{../figures/zebrafish}
\newcommand{\figdir}{../figures/zebrafish}

% No first-line paragraph indent; add a small inter-paragraph skip so the
% paragraph boundary stays visually obvious.
\setlength{\parindent}{0pt}
\setlength{\parskip}{0.5em plus 0.1em minus 0.05em}

\title{\textbf{Connectome-constrained self-motion integration in the larval
zebrafish hindbrain}}
\author{}
\date{}

\begin{document}
\maketitle
\begin{abstract}
How much computation can a circuit with a measured connectome support?
We test this on a vertebrate
self-motion integrator at single-neuron resolution. This integrator is the GABAergic
\rpi{} population of the larval-zebrafish anterior hindbrain and the
afferent populations the \emph{fish2} EM inventory identifies as
carrying its angular and translational drives. From the
\emph{neuprint-fish2} reconstruction we build sign-locked RNN and
message-passing GNN models whose recurrent connectivity is fixed to the
connectome under an all-inhibitory Dale convention, and train the
\emph{same} wiring on a family of progressively harder integration
tasks (heading, heading $+$ forward distance, 2-D path integration,
and a proprioceptive sensory-vs-efference gain mismatch) to ask what
the wiring \emph{can} compute.
\end{abstract}

\paragraph{From mechanism to geometry.}
Low-dimensional population structure is often analysed directly from neural
recordings, where rings, manifolds, and other geometric objects are extracted
as compact descriptions of activity. Such analyses typically focus on the
directions of largest variance, both because they are easiest to estimate and
because they are assumed to carry the dominant computations. Here we approach
the problem from the opposite direction. We begin with a measured connectome,
construct a dynamical model constrained by that wiring, train it to perform a
defined self-motion computation, and only then analyse the resulting population
activity. The geometric structures reported below are therefore not assumptions
and not purely descriptive summaries of data. They are consequences of a
specific recurrent computation implemented on a specific connectome.

This perspective is particularly important because the translation signal
recovered by the model occupies a low-variance population direction. Although it
is readily decoded from the full recurrent state, it is largely absent from the
dominant principal components. A purely variance-ranked analysis of activity
would therefore underestimate or entirely miss this component of the
computation. The geometry reported below should therefore be interpreted not
simply as a description of neural activity, but as a window onto which aspects
of the computation are visible in the dominant activity modes and which are
hidden in weaker yet behaviourally relevant population directions.

The translation estimate provides a cautionary example for variance-based
analyses of population activity. In the model, translation is readily recovered
because the computational variables are known: the task explicitly defines the
quantities of interest ($\theta$, $d$, and $(x,y)$), allowing us to test
directly whether they can be decoded from the recurrent state. In a recording,
by contrast, the latent variables are typically unknown and are often inferred
from the dominant activity modes themselves. Our results show that these two
approaches need not agree. Although translation is strongly decodable from the
full population, it contributes little to the total variance and is largely
absent from the leading principal components. A purely variance-ranked analysis
would therefore give little indication that such a variable exists, even though
it may be behaviourally important. The point is not that translation is
impossible to recover from neural data, but that identifying it requires either
an appropriate behavioural target or a mechanistic model that suggests where to
look.

\section{Anatomy, functions and connectome}
\label{sec:anatomy}

\paragraph{3-D anatomy.}
\Cref{fig:zhd_3d} shows the circuit cell families rendered in
3-D and colour-coded by the same seven-way functional partition used
in the circuit diagram: the recurrent ring substrate
(\emph{dIPN} $=$ \emph{IPNd*} $+$ \emph{IPNds*}), the \emph{IPN12}
pool, and the \emph{IPN-core} families \emph{IPN28/29/31--36}; the
three afferent populations that
the velocity gate routes the angular and translational drives through
(\emph{ARTR} $=$ \emph{RIPN01/02/03\_a/03\_b}, \emph{pt-IPN1}, and
\emph{motor\_efferent} $=$ \emph{RIPN11/12\_a/12\_c}); and a residual
``other'' category for sub-types not routed through the gate
(\emph{RIPN05/12\_b}). Spatially, the
\emph{dIPN} substrate forms the compact midline neuropil;
\emph{RIPN}-derived afferents descend from the habenula (a paired
diencephalic nucleus relaying limbic and modulatory signals to the
hindbrain) and \emph{pt-IPN}-derived afferents sweep in from the
lateral pretectum (a midbrain visual-processing region carrying optic-
flow and water-flow signals).

\begin{figure}[H]
  \centering
  \includegraphics[width=\textwidth]{\figdir/fig_1_overview.png}
  \caption{\textbf{Zebrafish HD / IPN circuit --- anatomy
  (\textbf{a}, \textbf{b}) and afferent partition (\textbf{c},
  \textbf{d}).} Seven-way colour code shared across panels:
  \emph{dIPN} (amber, $n=408$; ring substrate),
  \emph{IPN12} (rose, $n=108$; recurrent), \emph{IPN-core}
  (teal, $n=184$; recurrent),
  \emph{ARTR} (blue, $n=76$; angular afferent), \emph{pt-IPN1}
  (orange, $n=51$; exteroceptive translation), \emph{motor\_efferent}
  (green, $n=65$; angular/heading efference), \emph{other}
  (grey, $n=25$).
  \textbf{(a)}~Full \emph{neuprint-fish2} skeletons against Habenula /
  dIPN / vIPN ROI silhouettes; dorsal view.
  \textbf{(b)}~Same view restricted to cell bodies (SWC soma nodes).
  Panels (a, b) show spatially disjoint compartments of the same
  neurons: the dense central region of (a) is the dIPN / iIPN
  neuropil (no cell bodies), the somata sit in the surrounding cortex.
  \textbf{(c)}~Velocity-gate partition on the $917$-cell connectome:
  $\omega$ routed exclusively to ARTR and the exteroceptive forward
  channel $v_{\mathrm{ext}}$ routed exclusively to pt-IPN1; rotation
  and translation paths are anatomically separated.
  \textbf{(d)}~Proprioception extension: the motor-efferent RIPN
  sub-types (RIPN11/12\_a/12\_c) carry an efference copy of the
  \emph{heading/turn} command and are driven by a separate companion
  angular channel $\omega_{\mathrm{proprio}}$ in parallel with ARTR (green box
  between the ARTR and pt-IPN1 boxes), modelling sensory plus
  efference-copy redundancy on the \emph{angular} drive.
  Recurrent network and decoder are identical between (c) and (d);
  only the gate-target taxonomy differs.}
  \label{fig:zhd_3d}
\end{figure}

\paragraph{Functional circuit and afferent cell types.}
\label{method:proprioception}
The angular (heading) drive reaches IPN through two complementary
afferent pathways (\cref{fig:zhd_3d}\textbf{d}). One is the \emph{sensory}
ARTR population. The other is a \emph{motor-efference} copy of the turn
command. The exteroceptive \emph{translation} drive enters separately
through \emph{pt-IPN1}. Each afferent type carries a distinct, identifiable
signal. \emph{ARTR}, the Anterior Rhombomere Turning Region, is a
long-time-constant population. It typically becomes active only after the
first turning event, then sustains its activity over a relatively long
window. It is built from two similar pairs. \emph{RIPN01} and
\emph{RIPN03\_a} share a connectivity pattern that differs only in its ring
target: \emph{RIPN01} projects to \emph{dIPN}, and \emph{RIPN03\_a} to
\emph{dsIPN} ($=$\emph{IPNds*}), a small region just posterior to dIPN.
\emph{RIPN02} and \emph{RIPN03\_b} form a second pair with the targets
swapped (\emph{RIPN03\_b}$\to$\emph{dIPN}, \emph{RIPN02}$\to$\emph{dsIPN}).
This pair reads out the \emph{hSAG} pathway carrying horizontal rotational
signals. These latter cells are not yet reliably driven in the present data
(their activity resembles noise), but their connectivity places them with
ARTR. \emph{pt-IPN1} integrates optic-flow and water-flow signals. It
divides into roughly four groups, each encoding a different direction. The
\emph{motor-efference} cells are the motor-efferent RIPN sub-types
\emph{RIPN11}, \emph{RIPN12\_a}, and \emph{RIPN12\_c} ($36+29$ cells per
side, glutamatergic). The \emph{fish2} inventory ties them to head
direction, not forward swim. They split by efferent source. \emph{RIPN11}
receives \emph{ascending} motor efferents together with hSAG, and is our
leading candidate for moving the heading bump. \emph{RIPN12\_a} and
\emph{RIPN12\_c} (best treated as a single group) receive \emph{descending}
motor efferents together with hSAG. In the \emph{proprioception variant} of the model, these motor-efference
cells become a second, independent angular gate target. The sensory angular
velocity $\omega$ routes to \emph{ARTR}. A separate proprioceptive efference
copy routes to \emph{motor\_efferent}. The forward drive $v_{\mathrm{fwd}}$
routes only to \emph{pt-IPN1}. Keeping the two angular streams on distinct
ports is what lets us later drive them apart and ask whether the circuit
tracks their mismatch (\cref{method:mismatch}). This rotation/translation
split of the afferent partition (\cref{fig:zhd_3d}) mirrors the functional
organisation \citet{feierstein2023} report in the zebrafish hindbrain. They
apply reduced-rank regression to whole-brain GCaMP6s imaging. They find
that behaviour-related population activity is low-dimensional, captured by
two features. The corresponding voxels segregate into a
vergent-movement/swimming (translational) cluster and two rotation clusters
that lie selectively in the left and right hemispheres.

\paragraph{Connectome.}
\Cref{fig:zhd_connectome_summary} depicts the nominal $917$-cell
circuit, its binary support $\mathrm{supp}(W^{\mathrm{con}})$ and the
signed $W^{\mathrm{con}}$. The matrix is sparse (density $0.037$,
$30\,851$ non-zero edges out of $917^2 = 840\,889$); edge weights are
the \emph{fish2} synapse contact areas, spectrally normalised to
radius $0.9$. The recurrent substrate (the \rpi{} cells
\emph{IPNd*}, \emph{IPNds*}, the \emph{IPN12} pool, and the
\emph{IPN-core} families) is taken GABAergic: its
outgoing weights are sign-flipped to inhibitory, giving an
all-inhibitory recurrent substrate ($E\!:\!I = 10\,406\!:\!20\,445$,
$33.7\%$ / $66.3\%$). \emph{RIPN*} and \emph{pt-IPN*} afferents have
no Dale annotation on the \emph{fish2} server and are kept positive
(excitatory).

\begin{figure}[H]
  \centering
  \includegraphics[width=\textwidth]{\figdir/fig_2_connectome.png}
  \caption{\textbf{Zebrafish HD / IPN connectome (nominal $917$-cell
  circuit).} Matrix size $N \times N = 917 \times 917$, $30{,}851$
  non-zero edges (density $=0.037$); edge weights are \emph{fish2}
  synapse contact areas, the recurrent ring Dale-flipped inhibitory so
  $E\!:\!I = 10{,}406\!:\!20{,}445$ ($33.7\%$ / $66.3\%$). Both panels
  share the row/column ordering: neurons are sorted by the seven-way
  functional partition introduced in \cref{fig:zhd_3d}
  --- \emph{dIPN} (amber, $n=408$), \emph{IPN12} (rose, $n=108$),
  \emph{IPN-core} (teal, $n=184$), \emph{other} (grey, $n=25$),
  \emph{ARTR} (blue, $n=76$), \emph{pt-IPN1} (orange, $n=51$),
  \emph{motor\_efferent} (green, $n=65$) --- so the same block appears
  at the same matrix coordinates across (a) and (b), and the colour
  code matches \cref{fig:zhd_3d}.
  \textbf{(a)} Support mask. Colour-coded strips on the row and column
  axes mark each partition block, and translucent column bands inside
  the matrix carry the same colour code.
  \textbf{(b)} Signed $W^{\mathrm{con}}$, $z$-scored (red excitatory,
  blue inhibitory).}
  \label{fig:zhd_connectome_summary}
\end{figure}

\section{Methods}
\label{sec:methods}

\textbf{Recurrent neural network.}\label{method:RNN}
The RNN is a discrete-time recurrent network for
$i\in\{1,\dots,N\}$ with $N=917$:
\begin{equation}
\tau\,\frac{d\hat h_i(t)}{dt}
  = -\hat h_i(t)
  + \sum_j \hat W^{\mathrm{rec}}_{ij}\,\sigma\!\big(\hat h_j(t)\big)
  + \sum_l \hat W^{\mathrm{in}}_{il}\,u_l(t)
  + \hat b_i,
\label{eq:zhd_dynamics}
\end{equation}
$\sigma$ is the sigmoid, $\tau=0.1\,\mathrm{s}$, $\Delta t=0.01\,\mathrm{s}$.
The recurrent matrix is sign-locked to the connectome template
$W^{\mathrm{con}}\in\R^{N\times N}$ (built from the
\emph{neuprint-fish2} reconstruction; outgoing columns of GABAergic
HD-ring cells \emph{IPNd*}, \emph{IPNds*} are negated, and the whole
matrix is spectrally normalised to radius $0.9$): the trainable
per-edge parameter is the magnitude $\hat S$, and
\begin{equation}
\hat W^{\mathrm{rec}}_{ij}
  = |\hat S_{ij}|\,\mathrm{sign}(W^{\mathrm{con}}_{ij}).
\label{eq:zhd_w_rec}
\end{equation}
\bigskip
\noindent\textbf{Input gate on $\hat W^{\mathrm{in}}$}%
\label{method:input_extended}
$\hat W^{\mathrm{in}}\in\R^{N\times n_{\mathrm{in}}}$ has one row per
neuron ($N=917$) and one column per input channel
($n_{\mathrm{in}}=\dim u(t)$). At every step it injects the external
drive $\hat W^{\mathrm{in}}u(t)\in\R^{N}$ into the dynamics, added to
the recurrent term $\hat W^{\mathrm{rec}}\sigma(\hat h)$ of
\cref{eq:zhd_dynamics}. The rule is that each input variable may touch
only the one afferent cell type that
\cref{fig:zhd_3d}\textbf{c}--\textbf{d} assigns to it; the three
afferent types are the only entry points, and the ring cells
(\emph{IPNd*}, \emph{IPNds*}, \emph{IPN12}) never receive external
input. Take the exteroceptive variant of \cref{fig:zhd_3d}\textbf{c}, with the
four-channel drive
\begin{equation*}
u(t)=\bigl(\,\omega,\ v_{\mathrm{fwd}},\
   \cos\theta_0\!\cdot\!\delta_{t_0},\
   \sin\theta_0\!\cdot\!\delta_{t_0}\,\bigr)^{\!\top},\qquad
\omega\to\mathrm{ARTR},\quad v_{\mathrm{fwd}}\to\text{pt-IPN1}.
\end{equation*}
Grouping the $N$ rows by population, the signal injected into the
network at time $t$ (the $+\hat W^{\mathrm{in}}u(t)$ term of
\cref{eq:zhd_dynamics}) is the matrix--vector product
{\small
\begin{equation}
\begin{gathered}
\hat W^{\mathrm{in}}\,u(t)=
\bordermatrix{%
 & \omega & v_{\mathrm{fwd}} & \cos\theta_0 & \sin\theta_0 \cr
\mathrm{ARTR}_L   & -\softplus(\tilde v_\omega^{L}) & 0 & * & * \cr
\mathrm{ARTR}_R   & +\softplus(\tilde v_\omega^{R}) & 0 & * & * \cr
\text{pt-IPN1}_L  & 0 & v_{\mathrm{fwd}}^{L} & * & * \cr
\text{pt-IPN1}_R  & 0 & v_{\mathrm{fwd}}^{R} & * & * \cr
\text{ring}       & 0 & 0 & * & * \cr
\text{other aff.} & 0 & 0 & * & * \cr}
\begin{pmatrix}\omega\\[1pt] v_{\mathrm{fwd}}\\[1pt]
   \cos\theta_0\,\delta_{t_0}\\[1pt] \sin\theta_0\,\delta_{t_0}\end{pmatrix}
\\[6pt]
=\ \omega\,\mathbf{w}_\omega+v_{\mathrm{fwd}}\,\mathbf{w}_{v}
  +\delta_{t_0}\bigl(\cos\theta_0\,\mathbf{c}+\sin\theta_0\,\mathbf{s}\bigr)
  \ \in\R^{N},
\end{gathered}
\label{eq:zhd_swim_gate}
\end{equation}
}
where ``$*$'' marks a free (unconstrained) entry and
$\mathbf{w}_\omega,\mathbf{w}_{v},\mathbf{c},\mathbf{s}\in\R^{N}$ are
the four columns of $\hat W^{\mathrm{in}}$. The three column kinds are:
the \emph{sign-locked turn column} $\mathbf{w}_\omega$
($\omega\!\to\!$ARTR), nonzero only on ARTR with its left half forced
inhibitory and right half excitatory via $\mp\softplus$ (the two
scalars $\tilde v_\omega^{L},\tilde v_\omega^{R}$ are the only learned
parameters), and the lock makes left/right turns drive the ring peak in
opposite directions (a same-sign gate would integrate heading equally
well but drive both hemispheres in phase, unlike the data); the
\emph{unlocked vigour column} $\mathbf{w}_{v}$
($v_{\mathrm{fwd}}\!\to\!$pt-IPN1), non-directional, keeping the pt-IPN1
restriction but dropping the lock (one free scalar per side, either
sign); and the \emph{free cue columns} $\mathbf{c},\mathbf{s}$
($\cos\theta_0,\sin\theta_0$), dense and fully free over all $N$ rows,
which $\delta_{t_0}$ switches on only at the first frame $t_0$ to plant
the starting heading bump into the ring before they fall silent.
The proprioception variant of \cref{fig:zhd_3d}\textbf{d} adds
a fifth channel,
\begin{equation*}
u(t)=\bigl(\,\omega,\ v_{\mathrm{fwd}},\ \omega_{\mathrm{proprio}},\
   \cos\theta_0\!\cdot\!\delta_{t_0},\
   \sin\theta_0\!\cdot\!\delta_{t_0}\,\bigr)^{\!\top},\qquad
\omega_{\mathrm{proprio}}\to\mathrm{motor\_efferent},
\end{equation*}
i.e.\ a \emph{second} sign-locked angular column routed to
motor\_efferent (RIPN11 + RIPN12\_a + RIPN12\_c, which the \emph{fish2}
inventory ties to head direction, carrying the turn command's efference
copy), while $v_{\mathrm{fwd}}$ still routes only to pt-IPN1. In the
first model version the two angular columns share a driver,
$\omega_{\mathrm{proprio}}(t)=\omega(t)$, each with its own pair of
scalars. The same rule scales to the full multi-modal driver of the
nine ZAPBench conditions (\cref{fig:zhd_allblocks}): every visual
channel (rotation/optic-flow velocity, whole-field luminance, L/R light
difference, random-dot motion) drives pt-IPN1, and every tail-EMG
self-motion channel drives RIPN, each gated by \cref{eq:zhd_swim_gate}.
Anatomy fixes \emph{which} port a modality uses, the sign-lock fixes its
L/R \emph{polarity}, and only the gain $\softplus(\tilde v_c^{L,R})$ is
free.

\bigskip
\noindent\textbf{Output readout on $\hat W^{\mathrm{out}}$}%
\label{method:output_readout}
$\hat W^{\mathrm{out}}$ reads from the HD ring (\emph{IPNd*},
\emph{IPNds*}, \emph{IPN12 pool}) and never from the afferent
populations. Its column dimension $n_{\mathrm{out}}$ is set by the
active sub-task (\cref{fig:zhd_swim_task_modes}):
\begin{equation}
\hat y(t) =
\begin{cases}
(\cos\hat\theta,\sin\hat\theta)                 & n_{\mathrm{out}}=2 \ \ (\text{heading only}),\\
(\cos\hat\theta,\sin\hat\theta,\hat d)          & n_{\mathrm{out}}=3 \ \ (\text{heading}+\text{forward distance}),\\
(\cos\hat\theta,\sin\hat\theta,\hat x,\hat y)   & n_{\mathrm{out}}=4 \ \ (\text{2D path integration}),
\end{cases}
\label{eq:zhd_readout}
\end{equation}
with $\hat y = \hat W^{\mathrm{out}}\,\sigma\!\big(\hat h(t)\big)+\hat b^{\mathrm{out}}$,
$\hat\theta=\operatorname{atan2}(\hat y_2,\hat y_1)$. The leaky-vs-cumulative choice
(\cref{eq:xi_leaky,eq:xy_pi}) changes only what $\hat d$ and
$(\hat x,\hat y)$ are supervised against, not the readout shape.
The learnable parameters of the model are
\begin{equation*}
\big\{\hat W^{\mathrm{rec}}_{ij},\;
\{\tilde v_c^{L},\tilde v_c^{R}\}_c,\;
\hat W^{\mathrm{in}}_{:,\theta_0},\,\hat b,\,
\hat W^{\mathrm{out}},\,\hat b^{\mathrm{out}}\big\};
\end{equation*}
the connectome template $W^{\mathrm{con}}$, the afferent
subpopulation membership, the neuron-type assignments, $\tau$ and
$\Delta t$ are given.

\bigskip
\noindent\textbf{Graph neural network model.}\label{method:GNN}
The GNN model replaces the dense recurrent operator
$\sum_j \hat W^{\mathrm{rec}}_{ij}\,\sigma(\hat h_j)$ of
\cref{eq:zhd_dynamics} with a message-passing graph neural network
(Gilmer \emph{et al.}\ 2017). Compared to the RNN model, it does \emph{not}
hard-code either the leak, the firing-rate nonlinearity, or the
synaptic activation function. The per-edge message that combines a
presynaptic activity with its postsynaptic effect are optimized through learned functions $f_\theta$ and $g_\phi$. Beyond removing the functional bias, the GNN also admits an explicit form
of neuronal heterogeneity through a per-neuron learnable embedding
$\mathbf a_i \in \R^d$ that enters both $f_\theta$ and $g_\phi$.
Same-type neurons share their architectural role but can specialise
via their embedding, so a single pair of learnable functions covers the $34$ fish2
cell types of \cref{sec:anatomy} and all $N = 917$ instance-level
distinctions without any further per-type parameter tying. Nodes correspond to neurons and edges to the non-zero entries of the zebrafish HD connectome template
$W^{\mathrm{con}}$, and the subthreshold activity evolves as
\begin{equation}
\tau\,\frac{d\hat h_i(t)}{dt}
  = f_\theta\!\Big(\hat h_i(t),\,\mathbf a_i,\,
       \sum_{j \in \mathcal N_i}
       \hat W_{ij}\,g_\phi\!\big(\hat h_j(t),\mathbf a_j\big)\Big)
  + \sum_l \hat W^{\mathrm{in}}_{il}\,u_l(t),
\label{eq:zhd_gnn_dynamics}
\end{equation}
with the same readout (\cref{eq:zhd_readout}), swim-gated encoder
(\cref{eq:zhd_swim_gate}), and dynamics constants ($\tau,\Delta t$)
as the RNN. The neighbour set is fixed by the connectome topology,
$\mathcal N_i = \{j : W^{\mathrm{con}}_{ij} \neq 0\}$. The node
update $f_\theta$ and edge message $g_\phi$ are three-layer
perceptrons (width 64, ReLU, linear output). $g_\phi$ maps the
presynaptic state $\hat h_j$ to a per-edge message conditioned on
$\mathbf a_j$ and weighted by the per-edge gain $\hat W_{ij}$;
$f_\theta$ then processes the postsynaptic state $\hat h_i$, the
embedding $\mathbf a_i$, and the aggregated synaptic drive. We
constrain $g_\phi$ to be monotonically increasing in $\hat h_j$ via
the prior $\partial g_\phi / \partial \hat h \ge 0$ (see
\cref{eq:zhd_gnn_loss}); combined with the sign-locked $\hat W_{ij}$
(\cref{eq:zhd_gnn_w}) this restores Dale's law at the per-edge level
without introducing a $\pm$-sign degeneracy. The leak, the
firing-rate non-linearity, and any departure from a linear synapse
must emerge from $(f_\theta, g_\phi, \mathbf a)$ themselves; the
monotonicity priors of the training objective (\cref{eq:zhd_gnn_loss})
provide the minimum inductive bias for $f_\theta, g_\phi$ to behave like a
neuron (postsynaptic drive increases with net excitatory input). The per-edge weight is
sign-locked to the connectome as in the RNN,
\begin{equation}
\hat W_{ij}
  = |W_{ij}|\,\mathrm{sign}(W^{\mathrm{con}}_{ij}),
\label{eq:zhd_gnn_w}
\end{equation}
where $W_{ij} \in \R$ is the trainable per-edge magnitude.
In the zebrafish circuit $\mathrm{sign}(W^{\mathrm{con}}_{ij})$ is
negative for every edge whose presynaptic neuron is in the dIPN
(\emph{IPNd*} or \emph{IPNds*}) and positive for every edge from
\emph{RIPN*} or \emph{pt-IPN*}; the GNN therefore inherits the
inhibitory dIPN directly from the sign pattern. The learnable parameters are
\begin{equation*}
\big\{\theta,\,\phi,\,\mathbf a,\,\hat{W}_{ij},\;
\{\tilde v_c^{L},\tilde v_c^{R}\}_c,\;
\hat W^{\mathrm{in}}_{:,\theta_0},\,\hat b,\,
\hat W^{\mathrm{out}},\,\hat b^{\mathrm{out}}\big\};
\end{equation*}
the connectome template $W^{\mathrm{con}}$, the afferent
subpopulation membership, the neuron-type assignments, $\tau$ and
$\Delta t$ are given. As with the RNN, heading recovery from cos/sin readouts is
fundamentally non-unique: many distinct $(\theta,\phi,\mathbf a,W)$
configurations produce identical readouts, so the GNN is trained under a
regularised objective. This objective combines the same connectome-prior terms as the RNN with
sparsity and monotonicity priors on $f_\theta$ and $g_\phi$, specified
together with the rest of the training procedure in \cref{eq:zhd_gnn_loss}.

\bigskip
\noindent\textbf{Swim-integration task.}%
\label{method:swim_task_modes}\label{method:xi_leaky}
The on-disk swim-integration dataset is a \emph{single superset} from
which several sub-tasks are projected at training time. Swim events are
sampled at a Poisson rate $r_{\mathrm{swim}}=0.5\,\mathrm{Hz}$ with
event duration $\Delta t_{\mathrm{swim}}=0.3\,$s, partitioned into four
behavioural categories with fractions $L:R:F:B=0.3:0.3:0.2:0.2$. Left
and right events drive an angular signed phase impulse, forward and
backward events drive a signed forward-velocity impulse. The input is a
$4$-channel stream $u\in\R^{B\times T\times 4}=[\,\omega,\,v_{\mathrm{fwd}},\,
\cos\theta_0\!\cdot\!\delta_{t,0},\,\sin\theta_0\!\cdot\!\delta_{t,0}\,]$:
left/right events drive $\omega$, forward/backward events drive
$v_{\mathrm{fwd}}$, and the initial heading $(\cos\theta_0,\sin\theta_0)$
enters as a single-frame impulse at $t=0$ (position needs no such cue, as it
starts at the origin). The target column layout depends on the sub-task: the
\emph{scalar} recipe stores $y=[\,\cos\theta,\,\sin\theta,\,d\,]$ with
$\theta(t)=\theta_0+\int_0^t\omega\,ds$ and $d(t)=\int_0^t v_{\mathrm{fwd}}\,ds$
($d(0)=0$), while the \emph{2D path-integration} recipe stores
$y=[\,\cos\theta,\,\sin\theta,\,x,\,y\,]$, projecting $v_{\mathrm{fwd}}$ through
the current heading and integrating each spatial axis independently
(\cref{eq:xy_pi}). The 2D recipe is the \emph{true} path-integration target: it
couples translation to heading, so the network must maintain $\theta$
internally to predict $(x,y)$. Throughout the paper we use two words for the
two integrated quantities the circuit reports: the \emph{heading} (the
integrated angular self-motion $\theta$) and the \emph{translation} (the
integrated forward self-motion: the $1$-D distance $d$ in the scalar
recipe, the $2$-D position $(x,y)$ in the path recipe). The default scalar
$d$ and the default 2D position $(x, y)$ are both
\emph{cumulative} integrators of their respective drives. With a non-zero
mean forward-swim drive these targets grow linearly with the
curriculum horizon, so the loss scales with $T$ and cross-snapshot
comparisons become awkward. A finite time constant $\tau>0$ replaces
the cumulative sum with the leaky recurrence
\begin{equation}
d(t+\Delta t) = \left(1-\frac{\Delta t}{\tau_d}\right)d(t)
    + v_{\mathrm{fwd}}(t)\,\Delta t,
\qquad d(0)=0,
\label{eq:xi_leaky}
\end{equation}
equivalent to $\dot d = -d/\tau_d + v_{\mathrm{fwd}}$. For the
2D-position target the same leaky law is applied independently to each
axis,
\begin{equation}
\frac{dx}{dt} = -\frac{x}{\tau_p} + v_{\mathrm{fwd}}\cos\theta,
\qquad
\frac{dy}{dt} = -\frac{y}{\tau_p} + v_{\mathrm{fwd}}\sin\theta,
\qquad x(0)=y(0)=0,
\label{eq:xy_pi}
\end{equation}
which keeps the trajectory bounded near the recent egocentric
displacement. Biologically this is the natural choice: real neural
integrators leak, and what the fish actually encodes is a
low-pass-filtered estimate of \emph{recent} forward swim vigour and
recent displacement, not absolute long-horizon position. The
production runs use $\tau_d=\tau_p=0.5\,$s, with the
cumulative-integrator recipes preserved as separate variants so the loss
and metric comparisons stay clean.

\bigskip
\noindent\textbf{Proprioceptive-gain mismatch task.}\label{method:mismatch}
The proprioception circuit (\cref{method:proprioception}) lets us pose a
sharper question than redundant sensory$+$efference integration: can the
circuit track a \emph{discrepancy} between its two angular streams? We
drive ARTR with the observed angular velocity $\omega$ and
motor\_efferent with a \emph{gain-modulated} proprioceptive copy
$\omega_{\mathrm{proprio}}(t) = g(t)\,\omega(t)$, where the gain $g(t)$
is a piecewise-constant process drawn uniformly in $[0,1.5]$ that steps
to a new value every ${\approx}2\,$s. The supervised third readout column is
no longer a forward distance but the integrated mismatch
\begin{equation}
m(t) = \int_0^t \big(\omega(s) - \omega_{\mathrm{proprio}}(s)\big)\,ds
     = \int_0^t \big(1 - g(s)\big)\,\omega(s)\,ds,
\label{eq:zhd_mismatch}
\end{equation}
the accumulated sensory-vs-proprioceptive heading discrepancy. This is zero
when proprioception is veridical ($g\equiv1$) and grows whenever the
two streams disagree. The input is the $5$-channel proprioception layout
$[\omega, v_{\mathrm{fwd}}, \omega_{\mathrm{proprio}},
\cos\theta_0\delta_{t_0}, \sin\theta_0\delta_{t_0}]$ ($n_{\mathrm{in}}=5$)
and the readout is $\hat y=(\cos\hat\theta,\sin\hat\theta,\hat m)$.

\begin{figure}[t!]
\centering
\includegraphics[width=\textwidth]{\figdir/fig_3_swim_task.png}
\caption{\textbf{Swim-integration task and its target recipes for the
same trial.} \textbf{(a)}~$4$-channel input:
$\omega$, $v_{\mathrm{ext}}$, and the initial-heading impulse
$(\cos\theta_0,\sin\theta_0)\,\delta_{t,0}$ (non-zero only at $t=0$;
shown as stem markers carrying the per-trial value).
\textbf{(b)}~Scalar-$d$ target: $\cos\theta$, $\sin\theta$, and
forward distance $d$; perfect cumulative integrator (green) vs leaky
$\tau_d=0.5\,$s (red dashed).
\textbf{(c)}~The same trial's 2D trajectory in spatial coordinates,
computed from the $\omega$ and $v_{\mathrm{ext}}$ inputs of (a):
cumulative path (green) and leaky path (red dashed, $\tau_p=0.5\,$s).
\textbf{(d)}~Proprioceptive-gain mismatch task on the same $\omega$
(\cref{method:mismatch}): a piecewise-constant gain $g(t)\in[0,1.5]$
(top) sets the proprioceptive copy $\omega_{\mathrm{proprio}}=g\,\omega$
(orange, middle) against the sensory $\omega$ (green); the supervised
scalar (bottom) is the integrated mismatch
$\int(\omega-\omega_{\mathrm{proprio}})\,dt$.}
\label{fig:zhd_swim_task_modes}
\end{figure}

\paragraph{Training.}

\bigskip
Training optimises a per-frame MSE
on the heading readout with two connectome-prior regularisers and
an $\ell_1$ sparsity prior on the recurrent magnitudes,
\begin{equation}
\mathcal{L}
= \tfrac{1}{BT}\sum_{b,t} \|\hat y_{b,t} - y_{b,t}\|_2^2
+ \lambda_{\mathrm{cos}}\,\mathcal{L}_{\mathrm{cos}}
+ \lambda_{\mathrm{nf}}\,\mathcal{L}_{\mathrm{nf}}
+ \lambda_{1}\,\|\hat{W}^{rec}_{ij} \|_1,
\label{eq:zhd_loss}
\end{equation}
with block losses computed over type-pair masks
$\mathcal{B}_p \subset \{1,\dots,N\}^2$ (one block per ordered
pre/post neuron-type pair $p \in \mathcal{P}$),
\begin{align}
\mathcal{L}_{\mathrm{cos}}
  &= \tfrac{1}{|\mathcal{P}|}\sum_{p \in \mathcal{P}}
     \!\left[1 - \frac{\langle \hat W^{\mathrm{rec}}_{\mathcal{B}_p},\,
         W^{\mathrm{con}}_{\mathcal{B}_p}\rangle}
       {\|\hat W^{\mathrm{rec}}_{\mathcal{B}_p}\|\,
        \|W^{\mathrm{con}}_{\mathcal{B}_p}\|}\right],
\label{eq:zhd_cos}\\
\mathcal{L}_{\mathrm{nf}}
  &= \tfrac{1}{|\mathcal{P}|}\sum_{p \in \mathcal{P}}
     \relu\!\big(\kappa -
       \overline{|\hat W^{\mathrm{rec}}_{\mathcal{B}_p}|}\big)^2.
\label{eq:zhd_nf}
\end{align}
The cosine-distance term $\mathcal{L}_{\mathrm{cos}}$ anchors each
block to the connectome template; the norm-floor term
$\mathcal{L}_{\mathrm{nf}}$ prevents any type-pair block from
collapsing to zero by penalising the squared shortfall of the block
mean absolute weight $\overline{|\hat W^{\mathrm{rec}}_{\mathcal{B}_p}|}$
below a fixed floor $\kappa$ (typical value $\kappa = 0.05$). Because
$\mathcal{L}_{\mathrm{cos}}$ is scale-invariant,
$\mathcal{L}_{\mathrm{nf}}$ is required to prevent a degenerate
solution in which the cosine prior is satisfied by an arbitrarily
small recurrent magnitude. The type-pair vocabulary $\mathcal{P}$ is
the cartesian product of the 34 fish2 cell types identified in
\cref{sec:anatomy}, restricted to pairs with at least one non-zero
entry in $W^{\mathrm{con}}$.

\bigskip
\noindent\textbf{GNN training objective.}
The message-passing GNN (\cref{method:GNN}) is trained on the same per-frame
MSE and the same connectome-prior regularisers, augmented with sparsity and
monotonicity priors on its learned functions. As with the RNN, heading
recovery from cos/sin readouts is fundamentally non-unique: many distinct
$(\theta,\phi,\mathbf a,W)$ configurations produce identical readouts. We
constrain the solution space with the connectome-prior regularisers
$\mathcal{L}_{\mathrm{cos}}$, $\mathcal{L}_{\mathrm{nf}}$
(\cref{eq:zhd_cos,eq:zhd_nf}, computed on the per-edge gain surface
$\hat W_{ij}$ with floor threshold $\kappa$), together with sparsity
and monotonicity priors on the MLPs:
\begin{align}
\mathcal{L} \;=\;&
  \tfrac{1}{BT}\sum_{b,t}\|\hat y_{b,t} - y_{b,t}\|_2^2
+ \lambda_{\mathrm{cos}}\mathcal{L}_{\mathrm{cos}}
+ \lambda_{\mathrm{nf}}\mathcal{L}_{\mathrm{nf}} \nonumber\\
& +\;\lambda_0\|\theta\|_1 + \lambda_1\|\phi\|_1
   + \lambda_2\|\theta\| + \lambda_3\|\phi\|
   \;+\;\gamma_0\|W\|_1 + \gamma_1\|W\|_2 \nonumber\\
& +\;\mu_0\!\left\|
       \relu\!\left(+\tfrac{\partial f_\theta(\hat h,\mathbf a,m)}
                          {\partial \hat h}\right)\right\|
   +\;\mu_1\!\left\|
       \relu\!\left(-\tfrac{\partial f_\theta(\hat h,\mathbf a,m)}
                          {\partial m}\right)\right\| \nonumber\\
& +\;\mu_2\!\left\|
       \relu\!\left(-\tfrac{\partial g_\phi(\hat h,\mathbf a)}
                          {\partial \hat h}\right)\right\|
   +\;\mu_3\,\big\|g_\phi(h_\star,\mathbf a)-h_\star\big\|,
\label{eq:zhd_gnn_loss}
\end{align}
where $m = \sum_{j \in \mathcal N_i}\hat W_{ij}\,
g_\phi(\hat h_j,\mathbf a_j)$ is the aggregated synaptic input. The
$\ell_1$ penalties $\lambda_0,\lambda_1$ on the MLP weights bias
$f_\theta, g_\phi$ toward sparser, simpler functions; the $\ell_2$
penalties $\lambda_2,\lambda_3$ prevent weight explosion. The
leak-monotonicity term $\mu_0$ penalises \emph{positive} slopes of
$f_\theta$ in the postsynaptic state $\hat h$, encoding the prior
that the drift contains a leak
($\partial f_\theta / \partial \hat h \le 0$); $\mu_1$ penalises
negative slopes of $f_\theta$ in the aggregated input $m$, encoding
the prior that postsynaptic drive increases with net excitatory
input. Together, $\mu_0$ and $\mu_1$ force the discovered drift to be
decreasing in self-voltage and increasing in synaptic input, the
two qualitative signatures of \cref{eq:zhd_dynamics}. The term $\mu_2$
enforces the analogous monotonicity prior on $g_\phi$ in the
presynaptic voltage. The normalisation term $\mu_3$ anchors $g_\phi$
at a reference subthreshold value $h_\star$, removing the
multiplicative degeneracy between $g_\phi$ and $\hat W$ in
\cref{eq:zhd_gnn_dynamics}. Hyperparameters
$\lambda_{\{0,1,2,3\}},\,\gamma_{\{0,1\}},\,\mu_{\{0,1,2,3\}}$ and the
MLP widths/depths are tuned with an agentic optimisation procedure,
with $r_{\mathrm{roll},1k}$ (the Pearson correlation between the decoded
and true heading over a $1\,000$-step deterministic swim-impulse rollout)
as the primary score.

\bigskip
\noindent\textbf{Calcium observation loss.}\label{method:obs}
The heading loss of \cref{eq:zhd_loss} constrains only the
two-dimensional readout $\hat y(t)$; the $N$-dimensional recurrent state
$\hat h(t)$ is left unobserved, so the task alone does not pin the
internal dynamics. To constrain the implementation we add an optional
term that matches the model's predicted calcium to \emph{recorded}
activity. A GCaMP indicator turns the model's per-neuron activity into a
fluorescence trace by causal convolution with a unit-area
difference-of-exponentials kernel,
\begin{equation}
k_{\mathrm{GCaMP}}(t) \;\propto\; e^{-t/\tau_d} - e^{-t/\tau_r},
  \quad t \ge 0,
\qquad
\hat c_i(t) \;=\; \big(k_{\mathrm{GCaMP}} * \hat h_i\big)(t),
\label{eq:zhd_gcamp}
\end{equation}
with nominal cytosolic GCaMP7f rise/decay constants
$(\tau_r,\tau_d) = (0.15,\,1.20)\,\mathrm{s}$ \citep{dana2019}. ZAPBench
imaged the slower nuclear \emph{H2B-GCaMP7f}, so \cref{eq:zhd_gcamp} if
anything \emph{under}-states the low-pass smearing of the real observable.
The kernel is normalised to unit area so the transform
preserves the mean, the convolution is evaluated on the model grid
$\Delta t$, and it is differentiable, so gradients flow from the
fluorescence back into $\hat W^{\mathrm{rec}}$, the afferent gate and
$\hat b$. Let $\mathcal{O}$ index the $|\mathcal{O}| = 481$ HD-circuit
neurons matched to a recorded ROI (below) and let $c^{\star}_i(t)$ be the
recorded $\Delta F/F$ of neuron $i$. Because the indicator gain and
baseline are not identifiable, the term is made invariant to a per-neuron
affine rescaling by $z$-scoring both signals over time before comparison,
\begin{equation}
\mathcal{L}_{\mathrm{obs}}
  = \frac{1}{|\mathcal{O}|\,B_c T}\sum_{i\in\mathcal{O}}\sum_{b,t}
    \big(z[\hat c_i]_{b,t} - z[c^{\star}_i]_{b,t}\big)^2 ,
\qquad
z[x]_{b,t} = \frac{x_{b,t} - \mu_{b,i}}{\sigma_{b,i}},
\label{eq:zhd_obs}
\end{equation}
where $\mu_{b,i},\sigma_{b,i}$ are the per-trial, per-neuron mean and
standard deviation over time. For unit-variance signals
$\mathcal{L}_{\mathrm{obs}} = 2\,(1-\rho)$ with $\rho$ the Pearson
correlation, so the term rewards temporal \emph{shape} agreement rather
than absolute scale. The full objective adds it to \cref{eq:zhd_loss}
with weight $\lambda_{\mathrm{obs}}$,
$\mathcal{L}_{\mathrm{total}} = \mathcal{L}
  + \lambda_{\mathrm{obs}}\,\mathcal{L}_{\mathrm{obs}}$;
setting $\lambda_{\mathrm{obs}} = 0$ (or supplying no calcium trials)
recovers the task-only model exactly.

\paragraph{Mutual-information estimators.}\label{method:MI}
We use mutual information to quantify how strongly each unit's
firing rate $r_i(t) = \sigma(\hat h_i(t))$ tracks the heading
$\theta(t)$ (or, where stated, the angular velocity $\omega(t)$ or
the swim-onset trace $s(t)$) over a held-out rollout. For a single
neuron we wrap $\theta$ to $(-\pi, \pi]$ via
$\theta \leftarrow \arg(e^{i\theta})$, discretise the wrapped angle
into $n_\theta = 32$ equal-width angular bins and $r$ into $n_r = 20$
equal-width bins, then compute the plug-in estimator in bits,
\begin{equation}
\hat I(r_i;\theta) = \sum_{b,c} \hat p(b,c)\,\log_2
\frac{\hat p(b,c)}{\hat p(b)\,\hat p(c)},
\label{eq:zhd_mi_plugin}
\end{equation}
where $\hat p$ is the empirical joint over the
$(\theta\!\!-\!\!\mathrm{bin}, r\!\!-\!\!\mathrm{bin})$ histogram. The
plug-in construction is used so that the per-neuron and
per-cell-type quantities of
\cref{sec:results_four_classes} are
comparable across cell types and sub-populations. For a group of $K$ neurons we
use the cross-validated lower bound
\begin{equation}
\hat I_{\mathrm{LB}}(\mathbf r_{1:K};\theta) =
H(\theta) - \tfrac{1}{|\mathcal D_{\mathrm{te}}|}
\sum_{(t,\theta)\in\mathcal D_{\mathrm{te}}}
-\log_2 q_\psi(\theta_b\mid \mathbf r(t)),
\label{eq:zhd_mi_joint_lb}
\end{equation}
where $q_\psi$ is an $L_2$-regularised multinomial logistic
regression with input $\mathbf r \in \mathbb R^K$ and output the
$n_\theta$ heading bins, trained on a stratified 5-fold split.

\iffalse  % --- Image-domain ROI sampling / dF/F0 paragraph commented out ---
\paragraph{Image-domain ROI sampling and $\Delta F/F_0$.}\label{method:dff}
To make the model's population dynamics comparable to two-photon
calcium imaging of the larval-zebrafish dIPN, we sample activity
\emph{from the rendered frames} of the dorsal-view animation, not
from the latent state $\hat h_i(t)$. Each frame is a
$3000 \times 4380$-pixel canvas in which each neuron's skeleton is
drawn with a per-segment alpha proportional to its firing rate; we
place four axis-aligned ROI rectangles $R_k$
($k \in \{\mathrm{dIPN}_L, \mathrm{dIPN}_R, \mathrm{dsIPN}_L,
\mathrm{dsIPN}_R\}$) over the dorsal-IPN neuropil (drawn on frame~0
of \emph{\figdir/3D\_voltage\_const/frame\_0000.png}) and
integrate the green channel of the canvas RGBA buffer:
\begin{equation}
I_k(t) \;=\; \frac{1}{|R_k|}\sum_{(x,y)\,\in\,R_k}
              \mathrm{img}_t[y, x, \mathrm{green}].
\label{eq:zhd_roi_intens}
\end{equation}
The four traces $I_k(t)$ are then converted to a fractional change
with respect to a post-init baseline, mirroring the Seelig / Kim /
Green calcium-imaging convention:
\begin{align}
F_0[k] &\;=\; \bigl\langle I_k(t)\bigr\rangle_{t \,\ge\,
                t_{\mathrm{init}}} + \varepsilon, \\
\Delta F/F_0\,(k,t) &\;=\;
   \frac{I_k(t) - F_0[k]}{F_0[k]},
\label{eq:zhd_dff}
\end{align}
with $t_{\mathrm{init}} = 0.5$~s (drops the network start-up
transient that would otherwise dominate the baseline) and
$\varepsilon = 10^{-6}$. The per-pair population-vector position
overlaid on the kinographs in red is the activity-weighted centre
of mass across each $L/R$ pair,
$\mathrm{pos}_p(t) = \big(I_{p,L}\,y_{p,L} + I_{p,R}\,y_{p,R}\big)/
                      \big(I_{p,L} + I_{p,R}\big)$,
i.e.\ the same Seelig/Hulse PVA reduced to two phase bins per pair.
Image-domain sampling inherits all the projection biases of the
microscope (depth-of-field, projection, downsample) and so the
kinographs are directly comparable to \petrucco{}~Fig.~2 without
having to invert the calcium kernel.
\fi  % --- end commented-out ROI / dF/F0 paragraph ---

\section{Results}
\label{sec:results}

\subsection{Trained heading integrator}
\label{sec:results_evolution}

We train the sign-locked RNN of \cref{method:RNN} on the nominal
$917$-cell circuit of \cref{fig:zhd_3d}\textbf{c} (angular drive
$\omega$ routed exclusively to ARTR, \emph{RIPN01/02/03\_a/03\_b})
on the
heading-only sub-task of \cref{fig:zhd_swim_task_modes} ($n_{\rm in}=3$
channels, $n_{\rm out}=2$ outputs).
\Cref{fig:zhd_evolution_artr_pt1} reports the trained network.
After training the heading code is correctly recovered. Across the
five held-out swim-integration test trials (panel \textbf{i}) the
decoded heading tracks ground truth with a mean per-trial circular
RMSE of $5.9^\circ \pm 2.5^\circ$ and Pearson $r=0.997 \pm 0.003$.
Constant-$\omega$ deterministic rollouts of $T=2000$ frames spanning
the full $\pm 180^\circ/\mathrm{s}$ range (panels \textbf{h},
\textbf{j}) integrate heading at a near-unit gain: the linear-fit slope
of decoded vs.\ true $\omega$ is $1.07$ (mean local gain
$1.04 \pm 0.05$; panel \textbf{h}), and the Pearson correlation on the
unwrapped trajectories is $\ge 0.999$ for every
sweep. The dIPN carries a single travelling
ring activity peak (panel \textbf{e}). The learned recurrent weights
$\hat W^{\mathrm{rec}}$ sit on the signed connectome support by
construction; the per-edge magnitude correlates with the connectome
template at Pearson $r=0.430$ (panel \textbf{c}), consistent with the
sign-locked re-scaling described in \cref{method:RNN}.

\begin{figure}[p]
  \centering
  \includegraphics[width=0.9\textwidth]{\figdir/fig_evolution_zebrafish_hd_si_ipn_917_v1_selfmotion_rotation.png}
  \caption{\textbf{Trained heading integrator on the ARTR / pt-IPN1
  variant of \cref{fig:zhd_3d}\textbf{c}.} Sign-locked RNN with
  $\hat W^{\mathrm{rec}} = |\hat S|\odot\mathrm{sign}(W^{\mathrm{con}})$
  on the nominal $917$-cell circuit ($917$ v1, inhibitory). Drive:
  $\omega$ routed exclusively to ARTR
  (\emph{RIPN01/02/03\_a/03\_b}). Task: heading only ($n_{\rm in}=3$,
  $n_{\rm out}=2$).
  \textbf{(a)} GT connectome $W^{\mathrm{con}}$, neuron-type blocks on
  both axes.
  \textbf{(b)} Learned $\hat W^{\mathrm{rec}}$ on the same scale.
  \textbf{(c)} Scatter of the learned per-edge weight $\hat W_{ij}$
  against $W^{\mathrm{con}}_{ij}$ over all connectome edges.
  \textbf{(d)} All $917$ neurons over a constant-$\omega =
  +60^\circ/\mathrm{s}$ rollout, sorted by cell type.
  \textbf{(e)} dIPN sub-population over the same rollout; firing rates
  $z$-scored per neuron.
  \textbf{(f)} Constant-$\omega$ rollout. Top: $\omega(t)$. Bottom:
  true heading (green) and decoded heading (black), wrapped. Title
  reports post-warmup Pearson $r$ and circular RMSE.
  \textbf{(g)} Same layout on a held-out swim-integration test trial.
  \textbf{(h)} Integration gain (Hulse-style scatter): measured slope
  of the unwrapped decoded heading vs.\ true $\omega$, on a
  $15^\circ/\mathrm{s}$ grid over $[-180,+180]$; the
  near-linear domain is shaded, dashed line $g = 1$.
  \textbf{(i)} Five held-out swim-integration test trials: true
  heading (green) and decoded heading (black); per-panel circular RMSE
  and Pearson $r$.
  \textbf{(j)} Five constant-$\omega$ deterministic sweeps
  ($T=2000$ frames): same overlay; per-panel circular RMSE and Pearson
  $r$ on the unwrapped trajectories.}
  \label{fig:zhd_evolution_artr_pt1}
\end{figure}

\label{sec:results_decoder}
In \cref{sec:results_evolution} the heading was read out by a linear map
onto $(\cos\theta,\sin\theta)$ and decoded with $\operatorname{atan2}$.
The readout is not a neutral observer: because the recurrent weights are
connectome-locked and only re-scaled, the \emph{form} of the decoder is
one of the few degrees of freedom left to the optimiser, and it can
shape the internal representation the network settles on. We therefore
re-train the same $917$-cell circuit on the same heading task with two
alternative heading decoders, holding the recurrent dynamics and the
afferent gate fixed: (i) a \emph{$64$-bin softmax} readout, in which the
heading is one of $64$ angular bins (a localist ring code) decoded by the
softmax-weighted circular mean ($n_{\mathrm{out}}=64$); and (ii) a
\emph{nonlinear MLP} readout replacing the linear map.

The $64$-bin decoder converges as well as the $(\cos\theta,\sin\theta)$
readout: heading is recovered across held-out trials and constant-$\omega$
sweeps, and the dIPN still carries a single travelling activity peak, now
read out directly as a moving band in bin space
(\cref{fig:zhd_evolution_kbins}, panel \textbf{h}). The nonlinear MLP
decoder also converges, but only once the optimisation is \emph{paced}.
Trained at equal rates, the flexible MLP readout collapses to the target
mean (chance) before the connectome-locked ring can form. This is a
chicken-and-egg a linear readout avoids, because its optimum already
requires a structured ring. Letting the recurrent core lead and the
decoder follow (a higher, gently decaying recurrent learning rate against a
slower readout rate) removes the imbalance, and the MLP then recovers
heading as cleanly as the other two readouts (per-trial $r_\theta=0.999$,
median circular RMSE $1.3^\circ$; \cref{fig:zhd_evolution_mlpdec}). A small
capacity sweep confirms this is not tuned to one architecture: five
variants spanning hidden width $8$--$128$ and depth $2$--$4$ all converge
($r_\theta\ge0.994$, RMSE $1.3$--$3.9^\circ$). The takeaway is
methodological: with the recurrent substrate fixed to the connectome, the
decoder is an active modelling choice that interacts with the dynamics
(a too-flexible readout can out-run the circuit and must be paced so the
integrator forms first), but none of the three readouts (cos/sin,
localist bin, or nonlinear MLP) changes the underlying solution, and the
ARTR diagnostic confirms the recurrent dynamics are qualitatively unchanged
across all three readouts (\cref{fig:zhd_artr_oscillation}).

\label{sec:results_mismatch}
The proprioception circuit (\cref{fig:zhd_3d}\textbf{d}) routes the
angular drive through two anatomically distinct afferent streams: the
sensory ARTR copy $\omega$ and a gain-modulated motor-efference copy
$\omega_{\mathrm{proprio}}=g(t)\,\omega$
(\cref{method:mismatch}; \cref{fig:zhd_swim_task_modes}\textbf{d}).
Trained to report the integrated mismatch
$m=\int(\omega-\omega_{\mathrm{proprio}})\,dt$ alongside heading, the same
sign-locked $917$-cell RNN recovers both. Unlike the heading, which is
encoded as the unit vector $(\cos\theta,\sin\theta)$ and decoded by
$\operatorname{atan2}$, the mismatch $m$ is an unbounded cumulative
quantity and is supervised directly as a scalar, a plain regression
head trained under the same per-frame MSE. Across held-out swim trials the
decoded heading and the decoded mismatch each track ground truth at
per-trial Pearson $r\approx0.999$, and the mismatch trace follows the
piecewise-constant gain regime-switches frame by frame
(\cref{fig:zhd_evolution_mismatch}, panel \textbf{g} and row \textbf{j}).
The circuit therefore does not merely echo one afferent stream: it
maintains a running \emph{difference} of the two, accumulating the
sensory-vs-proprioceptive heading discrepancy, a quantity that is zero
only when proprioception is veridical. This is the connectome-locked
analogue of a forward-model error signal, and it shows the dual-afferent
wiring can support gain-mismatch detection, not just redundant
integration.


\subsection{Trained heading + distance integrator}
\label{sec:results_variants}

We do not know a priori which integration the larval circuit actually
performs: whether it maintains a scalar forward distance or a full
2-D position, and whether either is accumulated losslessly or leaked.
Rather than commit, we train all four combinations,
(distance vs.\ 2-D path) $\times$ (cumulative vs.\ leaky), on
the same connectome-locked circuit and compare. This brackets the
biologically plausible options: the cumulative target is the idealised
path integrator, the leaky target the bounded \emph{recent-displacement}
estimate a real leaky neural integrator would hold.

The same $917$-cell circuit and sign-locked RNN trains successfully
on every sub-task projection of the swim-integration superset
(\cref{fig:zhd_swim_task_modes}). Each variant changes only the
target column layout ($\hat W^{\mathrm{out}}$, \cref{method:output_readout})
and, where applicable, the leaky-integrator time constant
$\tau_d=\tau_p=0.5\,$s; the gate, connectome and recurrent dynamics
are unchanged across all five runs. The quantitative
trajectory-precision metrics below are aggregated from held-out test
trials on each trained checkpoint, with constant-$\omega$ and
constant-$v_{\mathrm{fwd}}$ deterministic rollouts at $T=2000$ frames
as the extrapolation probe.
\Cref{tab:zhd_subtask_metrics} summarises the headline numbers; full
panel layout in each figure matches \cref{fig:zhd_evolution_artr_pt1}.

\begin{table}[H]
\centering
\small
\resizebox{\textwidth}{!}{%
\begin{tabular}{l c c c c c c}
\toprule
 & \multicolumn{3}{c}{\textit{rollout accuracy} (Pearson $r$)} & \multicolumn{3}{c}{\textit{precision horizon} (60\,s naturalistic drive)} \\
\cmidrule(lr){2-4}\cmidrule(lr){5-7}
Variant & $r_\theta$ & $r_d$ & $\bar r_{xy}$ & $\tau_\theta$ (s) & $\tau_d$ (s) & $|g_\theta\!-\!1|$ \\
\midrule
\multicolumn{7}{@{}l}{\textit{Sign-locked RNN}} \\
Heading                      & $0.998 \pm 0.023$ & ---               & ---               & $\ge 60$ & ---      & $0.010$ \\
Heading $+$ distance (cumulative)
                             & $0.995 \pm 0.035$ & $0.989 \pm 0.053$ & ---               & $20.7$   & $5.7$    & $0.026$ \\
Heading $+$ distance (leaky) & $0.997 \pm 0.033$ & $0.997 \pm 0.027$ & ---               & $21.6$   & $\ge 60$ & $0.019$ \\
2-D path (cumulative)        & ---               & ---               & $0.760 \pm 0.268$ & $4.4$    & $4.4$    & $0.58$  \\
2-D path (leaky)             & ---               & ---               & $0.945 \pm 0.152$ & $18.5$   & $5.4$    & $0.024$ \\
\midrule
\multicolumn{7}{@{}l}{\textit{Message-passing GNN}} \\
Heading                      & $0.997 \pm 0.008$ & ---               & ---               & $12.7$   & ---      & $0.031$ \\
Heading $+$ distance (cumulative)
                             & $0.993 \pm 0.030$ & $0.991 \pm 0.034$ & ---               & $7.9$    & $6.6$    & $0.051$ \\
Heading $+$ distance (leaky) & $0.994 \pm 0.025$ & $0.998 \pm 0.003$ & ---               & $8.9$    & $38.2$   & $0.043$ \\
2-D path (cumulative)        & ---               & ---               & $0.837 \pm 0.209$ & $4.1$    & $3.8$    & $0.60$  \\
2-D path (leaky)             & ---               & ---               & $0.945 \pm 0.155$ & $4.0$    & $4.4$    & $0.061$ \\
Heading (cluster-move emb.)  & $0.997 \pm 0.011$ & ---               & ---               & $30.4$   & ---      & $0.012$ \\
Heading (frozen emb.)        & $0.996 \pm 0.011$ & ---               & ---               & $9.5$    & ---      & $0.051$ \\
\bottomrule
\end{tabular}}
\caption{Per-trial Pearson $r$ (mean $\pm$ SD on held-out test trials). RNN rows: across all $n=10\,000$ trials. GNN rows:
on a $n=512$ trial sample. $r_\theta$: correlation of unwrapped decoded
vs.\ ground-truth heading after a $10$-frame warmup. $r_d$: correlation
of decoded vs.\ ground-truth forward distance $d$. $\bar r_{xy}$: mean
of per-axis Pearson on $(x,y)$. Precision horizon, measured on a $60\,$s rollout
driven by a naturalistic Ornstein--Uhlenbeck $\omega(t)$ and $v_{\rm
fwd}(t)$ (mean over $4$ seeds). $\tau_\theta$: time the decoded heading
stays within $15^\circ$ of the integrated truth ($\ge 60$ = never crosses
within the rollout). $\tau_d$: same for displacement, with a \emph{relative} tolerance --- time
the decoded distance (scalar tasks) or $2$-D position stays within $20\%$
of the true displacement (evaluated once the true displacement exceeds
$0.5$ units, skipping the near-zero start). A relative band is used because
the cumulative targets are unbounded: only the leaky scalar-distance
integrator holds to $20\%$ for the whole rollout, whereas the cumulative
and $2$-D codes lose it within a few seconds (the cumulative ones partly
because they leave their trained range).
$|g_\theta\!-\!1|$: heading integration-gain error (slope of decoded vs.\
true unwrapped heading minus $1$), the $\omega$-independent driver of
drift. The message-passing GNN (\cref{method:GNN}) replaces the dense
recurrent operator with per-edge $g_\phi$ and per-neuron $f_\theta$ MLPs
over a learnable per-neuron embedding on the same connectome, gate, and
heads as the RNN. The last two GNN rows are embedding ablations on the
heading task --- \emph{cluster-move} (the embedding stays clustered by cell
type but the clusters may drift during training) and \emph{frozen} (the
embedding is held at its cell-type initialisation).}
\label{tab:zhd_subtask_metrics}
\end{table}

\Cref{fig:zhd_evolution_both,fig:zhd_evolution_both_leaky} train on
the joint heading-plus-forward-distance task
($n_{\rm in}=4$, $n_{\rm out}=3$, target $(\cos\theta,\sin\theta,d)$).
Per-frame test R² is $\ge 0.9998$ on both heading and distance for
both variants. The constant-$\omega$ trajectory metrics over four
sweeps give a mean heading RMSE of $33.1^\circ$ (slope ratio $0.97$)
for the cumulative variant and $8.1^\circ$ (slope ratio $0.99$) for the
leaky variant: \emph{the leaky-target supervision actually produces
a sharper compass} than the cumulative-target version of the joint task
(and is comparable to the heading-only model of
\cref{fig:zhd_evolution_artr_pt1}). The distance head behaves
like the translation-only models: average $d$-RMSE over the five
constant-$v_{\mathrm{fwd}}$ sweeps is $10.9$ (cumulative) and $13.9$
(leaky), with slope-ratio $\dot{\hat d}/\dot d$ of $0.14$ (cumulative)
/ $0.003$ (leaky), confirming that the constant-input extrapolation
problem of the translation-only run is task-, not input-, dependent.

\Cref{fig:zhd_evolution_position2d,fig:zhd_evolution_position2d_leaky}
train on the full 2-D path task ($n_{\rm in}=4$, $n_{\rm out}=4$,
target $(\cos\theta,\sin\theta,x,y)$ with
$\dot x = v_{\mathrm{fwd}}\cos\theta$,
$\dot y = v_{\mathrm{fwd}}\sin\theta$). Heading is still recovered at
test R² $\ge 0.994$ on the disk-stored target. The 2-D-position head
is the harder regression: averaged over five $(\omega, v_{\mathrm{fwd}})$
sweeps at $v_{\mathrm{fwd}}=+1$, the per-frame Euclidean RMSE between
decoded $(\hat x,\hat y)$ and true $(x,y)$ is $2.10$ for the cumulative
variant and $0.98$ for the leaky variant: the leaky target
\emph{halves} the trajectory error. Mean endpoint error
$\|(\hat x,\hat y)_T - (x,y)_T\|$ similarly drops from $3.72$ to
$1.78$. The path-length ratio
$\sum\|\Delta(\hat x,\hat y)\|/\sum\|\Delta(x,y)\|$ is highly
variable across sweeps (cumulative $0.23$--$1.85$, leaky $0.28$--$0.98$),
indicating that the absolute arc-length is unreliable on deterministic
constant-drive probes even when the per-frame error is small. Panels
(f)--(g) of these figures show the decoded $(x,y)$ trajectory as a
2-D path instead of the heading time-trace.


\begin{figure}[p]
  \centering
  \includegraphics[width=0.9\textwidth]{\figdir/fig_evolution_zebrafish_hd_si_ipn_917_v1_position_2d_leaky.png}
  \caption{\textbf{Trained 2-D path integrator, leaky target}
  ($\tau_p=0.5\,$s). Panels follow \cref{fig:zhd_evolution_artr_pt1};
  panels (f)--(g) show the decoded path $(x,y)$ in spatial coordinates on
  the constant-$\omega$ rollout and on a held-out swim trial. The
  cumulative-target companion and the other model dashboards are in
  \cref{app:dashboards}.}
  \label{fig:zhd_evolution_position2d_leaky}
\end{figure}

\label{res:gnn_subtasks}
The same connectome-locked circuit also trains under the message-passing
GNN of \cref{method:GNN}, which replaces the dense recurrent operator
$\sum_j\hat W^{\mathrm{rec}}_{ij}\sigma(\hat h_j)$ with per-edge $g_\phi$
and per-neuron $f_\theta$ MLPs over a learnable per-neuron embedding
$\mathbf a_i$. The leak, the firing-rate nonlinearity, and the synaptic
activation are no longer hard-coded; they must emerge from
$(f_\theta, g_\phi, \mathbf a)$. We train the GNN on three of the
swim-integration sub-tasks (heading-only, 2-D path with the
cumulative target, and 2-D path with the leaky target) on the same
on-disk dataset, the same anatomy-routed afferent gate, and the same
heading and $(x, y)$ readout heads as the sign-locked RNN. The bottom
block of \cref{tab:zhd_subtask_metrics} reports the per-trial Pearson
$r$. On heading the GNN reaches $r_\theta = 0.997 \pm 0.008$, within
the SD of the RNN baseline ($0.998 \pm 0.023$). On the harder 2-D path
the GNN matches the RNN on the leaky target
($\bar r_{xy} = 0.945 \pm 0.155$ vs.\ $0.945 \pm 0.152$) and slightly
exceeds it on the cumulative target
($0.837 \pm 0.209$ vs.\ $0.760 \pm 0.268$). Lifting the hard-coded
leak, firing-rate nonlinearity, and synaptic activation does not
collapse the task solution: the circuit reaches the same
per-frame accuracy under a substantially less constrained dynamical
parameterisation, consistent with the integration being a property of
the connectome and the afferent partition rather than of a particular
choice of single-cell dynamics.
\Cref{fig:zhd_gnn_dashboard} shows the trained heading-only GNN in the
same panel layout as \cref{fig:zhd_evolution_artr_pt1}, with an extra
row reporting the three quantities the GNN adds over the sign-locked
RNN: the per-neuron embedding and the two learned functions (the
per-edge signalling function and the node update) that replace the
fixed leak and firing-rate nonlinearity. Trained on the joint targets the
GNN also reproduces the recurrent code: its pooled mutual information shows
the same distributed structure as the RNN (single neurons carry little,
the whole pool reaches the ceiling) across all four joint-target
variants (\cref{fig:zhd_mi_partition_gnn}, the companion of
\cref{fig:zhd_mi_partition}).


\subsection{Population structure of the distributed self-motion code}
\label{sec:results_four_classes}

The previous sections established that the same connectome-constrained circuit
can solve a family of self-motion integration tasks, from heading estimation to
full $2$-D path integration. The remaining question is how these variables are
represented within the recurrent population. In recent years, low-dimensional
population geometry has become a common language for describing neural activity,
with rings, manifolds, and other geometric structures routinely extracted from
recordings. Such descriptions are often informative, but by themselves they
remain descriptive: a geometric object observed in data does not explain why it
exists or what computation generates it. Here we approach the problem from the
opposite direction. We begin with a measured connectome, construct a dynamical
model constrained by that wiring, train it to perform a defined self-motion
computation, and only then analyse the resulting population activity. The
geometric structure reported below is therefore not an assumption and not merely
a compact summary of data. It is a consequence of a specific computation
implemented by a specific recurrent circuit.

This distinction is particularly important because behaviourally relevant
variables need not coincide with the dominant directions of population variance.
Previous work has shown that neural information can be distributed across
low-variance dimensions of activity, rather than being confined to the leading
principal components. Here we find an extreme example: heading dominates the
activity variance and forms the familiar ring-attractor geometry, whereas
translation occupies a distinct population direction that is strongly decodable
from the full recurrent state yet largely absent from the dominant principal
components. A variance-ranked analysis alone would therefore capture the heading
manifold while substantially underestimating the translation representation. The
goal of this section is to characterise how heading and translation are
distributed across neurons, cell types, and population dimensions, and to relate
that organisation back to the underlying connectome-constrained computation.

Recall the two integrated variables the circuit reports
(\cref{method:swim_task_modes}): the \emph{heading} (integrated angular
self-motion) and the \emph{translation} (integrated forward self-motion: a
$1$-D distance or a $2$-D position). We characterise the recurrent code in three
steps: which variable each neuron carries, the population geometry that supports
these representations, and how that geometry compares to the reduced-rank
population modes reported by \citet{feierstein2023} in the same hindbrain
region.

For the four models trained to
report heading together with a translation target (a $1$-D distance or a $2$-D
path, each cumulative or leaky) we ask where each variable lives. For every
recurrent neuron we measure the mutual information (\cref{method:MI}) its
activity carries about heading and about translation (the radial distance
for the $2$-D models), then decode the same two quantities from progressively
larger \emph{groups}: each fish2 cell type, each anatomical family, and the
whole $700$-cell GABAergic pool (the afferents, whose drive-locked MI is
trivially high, are excluded). \Cref{fig:zhd_mi_partition}, and its
message-passing-GNN companion \cref{fig:zhd_mi_partition_gnn}, plot all of
these in one heading-versus-translation plane, dot area growing with group
size.

\begin{figure}[H]
  \centering
  \includegraphics[width=0.8\textwidth]{\figdir/fig_zebrafish_mi_partition.png}
  \caption{\textbf{Pooled mutual information of the recurrent code: more
  neurons, more information.} The four task models, one per panel:
  \textbf{(a)}~heading\,$+$\,distance (cumulative),
  \textbf{(b)}~heading\,$+$\,distance (leaky), \textbf{(c)}~2-D path
  (cumulative), \textbf{(d)}~2-D path (leaky), all on the nominal
  $917$-cell circuit. Each panel plots mutual information about heading
  ($x$) against translation ($y$), in bits, on a common $0$--$4$ scale.
  Every marker is the joint MI of a \emph{group} of recurrent neurons
  (cross-validated decoder lower bound, \cref{eq:zhd_mi_joint_lb}; for the
  2-D models the translation target is the joint $(x,y)$ position), with
  \textbf{dot area growing with group size}: the faint cloud is the single
  neurons (per-neuron plug-in MI, $n=1$, coloured by family), then one dot
  per fish2 cell type and the whole recurrent pool (black, $n=700$). Family
  colours: \emph{dIPN} amber, \emph{dsIPN} teal, \emph{IPN12} rose,
  \emph{IPN-core} purple. The single-neuron plug-in MI and the group
  decoder lower bound are different estimators (not comparable
  point-by-point) but consistent across panels. Light grey diagonals mark
  constant \emph{total} information --- heading MI $+$ translation MI ---
  at each integer bit (labelled in \textbf{b}); a group's distance up the
  diagonals is its total joint code. The family-colour key is in panel
  \textbf{a}, the dot-size key in panel \textbf{d}.}
  \label{fig:zhd_mi_partition}
\end{figure}

Heading and translation are represented very differently. Heading is encoded redundantly at every scale, from single neurons to the entire recurrent population. Many individual neurons are strongly heading-tuned, and any sufficiently large group decodes heading near the $\approx 3.7$-bit ceiling. Translation depends strongly on the task. For the 1-D forward-distance target (\cref{fig:zhd_mi_partition}\textbf{a},\textbf{b}), some individual neurons are distance-tuned and the per-neuron cloud extends above $1$ bit on the translation axis. For the 2-D path-integration target (\cref{fig:zhd_mi_partition}\textbf{c},\textbf{d}), however, essentially no single neuron carries substantial position information (per-neuron translation information $\approx 0$), yet the recurrent population as a whole decodes position near the information ceiling ($\approx 3.4$ bits), consistent with the path-integration accuracy of the network ($\bar r_{xy}=0.95$, \cref{tab:zhd_subtask_metrics}).\footnote{This agreement is a consistency check rather than an exact identity. For jointly Gaussian variables with correlation $\rho$, $I=-\tfrac12\log_2(1-\rho^2)$. Two approximately independent position coordinates each decoded at $\rho\approx0.95$ yield $\approx3.4$ bits, while the heading coordinate decoded at $\rho\approx0.997$ yields $\approx3.7$ bits.} The 2-D position estimate is therefore represented as a genuinely distributed population code: it is readily recoverable from the population yet largely absent at the level of individual neurons.

This division is computational rather than anatomical. No recurrent family occupies a distinct region of the heading-versus-translation information plane of \cref{fig:zhd_mi_partition}. All four GABAergic families (\emph{dIPN}, \emph{dsIPN}, \emph{IPN12}, \emph{IPN-core}) contribute to both variables with only modest biases, despite the fact that heading and translation enter through anatomically distinct afferent pathways (\cref{fig:zhd_3d}\textbf{c}). The recurrent computation therefore mixes the two streams rather than preserving the input partition. The code is also highly redundant. Pooled information grows only logarithmically with group size ($I_{\mathrm{group}}\approx a+b\log_{10}N$, with $b\approx0.80$ bits/decade for heading and $0.93$ bits/decade for translation; \cref{fig:zhd_mi_grouped}), indicating strong overlap between neurons. The full $700$-cell recurrent pool already lies close to saturation, so adding additional neurons contributes little extra information. The message-passing GNN reproduces the same qualitative organisation (\cref{fig:zhd_mi_partition_gnn}).


\begin{figure}[t!]
  \centering
  \includegraphics[width=\textwidth]{\figdir/fig_zebrafish_distributed_code.png}
  \caption{\textbf{Heading forms a dominant ring manifold, whereas translation
  emerges as a distributed population code occupying a different population
  direction.} Recurrent firing rates over held-out swim rollouts; panels
  \textbf{a}--\textbf{e} use the $2$-D-path model, \textbf{f} all four
  joint-target models, which share one neuron basis.
  \textbf{(a)}~The population projected onto its two heading axes, coloured by
  true heading: a clean ring attractor (participation ratio and cross-validated
  heading decode annotated).
  \textbf{(b)}~The same population in a $3$-D frame built from the two heading
  axes and the translation axis orthogonal to the heading plane, coloured by
  true distance: the translation axis lifts almost vertically out of the
  heading ring (their principal angle annotated), i.e.\ translation occupies a
  distinct population direction.
  \textbf{(c)}~Translation information (cross-validated decoder lower bound)
  against group size, from single neurons through cell types and anatomical
  families to the whole recurrent pool (colours mark families; dashed line, the
  logarithmic trend). Single neurons carry little $2$-D position information; the
  whole pool carries it near the entropy ceiling, an order of magnitude more.
  \textbf{(d)}~Translation decode $R^2$ against the number of \emph{top-variance}
  principal components retained (curves, the two $2$-D models; star, the
  variance-matched top-$2$ components) versus the task-tuned readout the network
  actually learns (dotted line). The estimate is rank-$2$-readable but oriented
  away from the directions of greatest activity: important does not mean
  high-variance.
  \textbf{(e)}~Reduced-rank regression after \citet{feierstein2023}: every
  neuron in the plane of its heading-feature and translation-feature loadings,
  coloured by population. The exteroceptive translation afferent loads on the
  translation axis and the two halves of the heading afferent split along
  $\pm$ heading --- three clusters as in their data, here \emph{built in} by the
  anatomy-routed input gate and L/R sign-lock --- while the recurrent pool (the
  emergent part) forms a mixed cloud through the centre.
  \textbf{(f)}~Cross-task geometry: the heading ring is reused across tasks
  (small mean principal angle) while the translation subspace rotates from the
  $1$-D distance axis to the $2$-D path plane (measured angles annotated).}
  \label{fig:zhd_distributed_code}
\end{figure}

Projecting the recurrent activity onto its principal axes reveals a strikingly low-dimensional population geometry: across all four task variants the participation ratio is only $3.2$--$4.7$ and $90\%$ of the variance is contained within $5$--$8$ dimensions, despite the recurrent population containing roughly $700$ neurons. The self-motion estimate therefore occupies a compact low-dimensional manifold. A targeted dimensionality reduction based on the task variables separates this manifold into two dominant components. A heading subspace carries most of the variance and forms a clean ring as heading sweeps through $360^\circ$ (\cref{fig:zhd_distributed_code}\textbf{a}). A second translation-related subspace lies nearly orthogonal to the heading ring (principal angle $84^\circ$--$89^\circ$ for the $1$-D-distance models and $68^\circ$--$86^\circ$ for the $2$-D-path models; the two position axes are themselves approximately orthogonal, $71^\circ$--$82^\circ$). Projected into a $3$-D frame spanned by the two heading axes and the translation axis, the population traces a ring in the heading plane while the translation direction lifts almost vertically out of it (\cref{fig:zhd_distributed_code}\textbf{b}), making the orthogonality directly visible.

This geometric separation explains the contrast between the single-neuron and population-level analyses. The heading estimate is visible at the level of individual neurons, whereas the $2$-D position estimate is primarily a population-level quantity: single neurons carry little information about the planar position, but pooling cell types, then anatomical families, then the whole recurrent pool lifts the decodable translation information by an order of magnitude, up towards the entropy ceiling (\cref{fig:zhd_distributed_code}\textbf{c}). A linear decoder restricted to the top two principal components recovers the $2$-D position poorly ($R^2\approx0.01$--$0.06$), whereas the task-trained readout learned by the network recovers it almost perfectly ($R^2\approx0.99$; \cref{fig:zhd_distributed_code}\textbf{d}). Thus the position estimate is encoded in a low-dimensional subspace that is not aligned with the directions of greatest activity variance. Translation is therefore readily decodable from the population yet difficult to recover from either single neurons or variance-ranked dimensionality reduction alone. A notable feature is that the heading geometry is conserved across tasks whereas the translation geometry is not. The heading ring remains largely unchanged between the distance and 2-D-path tasks, while the translation subspace rotates substantially (\cref{fig:zhd_distributed_code}\textbf{f}). The recurrent circuit therefore appears to reuse a common heading scaffold while adapting the representation of translation to the demands of the task.


Because the four models share one neuron basis, their subspaces can be compared
directly (\cref{fig:zhd_distributed_code}\textbf{f}). The heading subspace is
conserved across all four tasks (cross-task principal angles $10$--$30^\circ$,
nearly identical between the cumulative and leaky variants), whereas the
translation subspace rotates as the task grows: the $1$-D-distance axis sits
$60$--$67^\circ$ from the $2$-D-path plane, and the two $2$-D subspaces differ
by ${\approx}50^\circ$ between cumulative and leaky. The circuit reuses a
single heading ring as a stable scaffold and re-orients a separate
translation subspace for each task. The message-passing GNN reproduces this
within-model geometry: low-dimensional ($\mathrm{PR}\approx5$--$11$),
near-orthogonal heading and translation axes ($83$--$87^\circ$), translation
invisible to the top two modes ($R^2\!\approx\!0.02$--$0.06$) yet
rank-2-readable ($R^2\!\approx\!0.998$). The geometry is therefore a property of
the connectome-locked computation, not of the single-cell parameterisation.

\Citet{feierstein2023} analysed this hindbrain with reduced-rank regression
and found that behaviour-related activity collapses onto an effectively
two-dimensional subspace: one vergence / translation-plus-swim mode and one
rotation mode (their term for the angular self-motion we call heading), with
neurons separating by their loadings into a translation cluster and two
left / right rotation clusters that segregate to opposite hemispheres. Running
the identical analysis on the model (population activity regressed on
time-shifted behavioural regressors grouped into a translation and a heading
family) recovers the same low-dimensional behaviour subspace, whose two
dominant features split cleanly into a \emph{heading} feature (loading on
angular velocity and heading) and a \emph{translation} feature (loading on
forward velocity, swim, and displacement), the model's two behavioural modes.
In the plane of these two feature loadings the regression finds the same clean
three-cluster, hemispherically-split organisation \citet{feierstein2023}
report: the exteroceptive translation afferent loads on the translation axis,
the heading afferent on the heading axis with its two halves carrying
opposite-signed angular-velocity weights, while the recurrent pool forms a
mixed cloud through the centre (\cref{fig:zhd_distributed_code}\textbf{e}).
This separation, however, is
\emph{built in}, not discovered: the input gate routes the angular drive only
to the heading afferent and the forward drive only to the translation
afferent, with the heading afferent's two halves sign-locked to opposite
polarity (\cref{method:input_extended}). The regression at those cells
therefore reads the imposed wiring back out, a positive control that the
analysis recovers the input organisation we built in, which is itself the
anatomical organisation \citet{feierstein2023} describe and which is precisely
why we route the gate that way. The \emph{emergent} result is what the
recurrent integrator does with these cleanly separated inputs: it does
\emph{not} inherit their separation. Its neurons are mixed, with
heading-minus-translation selectivity near zero, the population-level echo
of the distributed code above. So where imaging of the (largely afferent /
sensorimotor) hindbrain reveals separate translation and rotation populations,
the connectome-locked recurrent computation \emph{downstream} of those inputs
distributes heading and translation across a shared, low-dimensional manifold
rather than segregating them. That mixing, not the imposed afferent split,
is the model's prediction, and the explicit geometric account of what
``distributed'' means here that \cref{out:manifold} called for.

\subsection{Task accuracy does not imply realistic activity}
\label{sec:results_task_vs_activity}
Solving the task is not the same as matching the brain. Driven by the
recorded stimulus, the task-only circuit decodes heading well yet produces
calcium quite unlike the recording. This is the core degeneracy that a
connectome plus a task loss leaves unresolved, since many circuits realise
the same computation with different activity. A calcium-observation loss
that pulls predicted activity toward the recording partially closes this
gap, and the fish's own behaviour bounds how far any such model can reach.
We develop these observation-matching experiments (the task-only
baseline, the behavioural covariates, the observation-loss sweep, and the
proprioceptive 2-D variant under calcium supervision) in
\cref{app:calcium}; the discussion and outlook below build on them.


\subsection{ARTR: opponent L/R gate and the origin of the oscillation}
\label{sec:results_artr}

The angular afferent pool (ARTR; \cref{tab:zhd_mapping}) is the anterior
rhombencephalic turning region of \citet{dunn2016artr}, defined biologically
by a slow ($\sim$10--20\,s) antiphase oscillation between its left and right
clusters that biases successive swim bouts. We ask three things of the
connectome-locked model: whether it reproduces the L/R opponency, what
generates the alternation, and whether the rhythm is intrinsic or inherited.
Every diagnostic is shown for the three heading decoders of
\cref{sec:results_decoder} together (the linear
$(\cos\theta,\sin\theta)$ readout, the $64$-bin localist readout, and the
nonlinear MLP), so any decoder dependence is visible at a glance.

The recorded L/R anti-correlation is not imposed by the heading task. Trained
unconstrained, the angular gate converges to a same-sign solution that
integrates heading equally well but drives both hemispheres in phase. Locking
the gate to opposite polarity (left $\le 0$, right $\ge 0$, magnitudes free)
restores opponency at no cost to accuracy; the non-directional channels
(forward vigour, luminance) carry no lock.Band-splitting the L/R population means at a $16$\,s period separates two
components (\cref{fig:zhd_artr_kino}). The slow band ($>16$\,s) is the
turn-block: antiphase and command-locked in both data ($r_{LR}=-0.57$) and
model ($-0.96$). The fast band ($<16$\,s) is the heading-bump carrier at the
rotation period $360/|\omega|$; it is L/R phase-locked in the model
($|r_{LR}|=0.90$) but incoherent and left-lateralised in the recording
($r_{LR}=+0.03$). An inherited carrier is necessarily L/R-locked, so the
recorded fast band is a graded mix: an inherited carrier plus a
cell-specific component at or below the imaging-SNR floor. The modelled ARTR
is antiphase and command-locked under all three decoders; the opponency is
strongest for the linear readout and weaker but preserved for the bin and MLP
readouts (whole-trace $r_{LR}=-0.93$, $-0.59$, $-0.53$), so the decoder leaves
a measurable imprint without changing the qualitative picture. 

An intrinsic oscillator has a period set by its own time constants
(invariant to the drive, and sustained without it), whereas an inherited one
tracks the drive and stops with it. Under a constant-$\omega$ sweep the ARTR
period tracks the bump period $360/|\omega|$ while its amplitude stays roughly
constant, for all three decoders (\cref{fig:zhd_artr_oscillation}): the
signature of a carrier slaved to the rotating heading bump, not a clock with
its own frequency. This is structural: the modelled ARTR is all-excitatory
with a single fast time constant, so it latches rather than oscillates. An
intrinsic rhythm would need recurrent left$\leftrightarrow$right inhibition
(for bistability) and a slow adaptation variable (to set the $\sim$10--20\,s
period), neither of which the connectome-locked circuit contains.

The decisive test removes the command: the ZAPBench session runs a rotation
block straight into darkness in the same fish. Driven identically (recorded
$\omega$, then zero input), the recorded ARTR keeps alternating in the dark
($|L-R|$ std $=0.72$) while the model is near-flat ($0.14$,
$\sim$5$\times$ smaller; \cref{fig:zhd_artr_dark}); a zero-input free-run from
a small perturbation decays to a fixed point for every decoder
(amplitude $<0.006$), confirming the model carries no self-sustained
oscillation. The recording carries an intrinsic dark-state rhythm
\citep{dunn2016artr} that the command-driven model lacks. This is a
prediction--data gap that names the missing ingredients (recurrent inhibition
$+$ slow adaptation) and a falsifiable target for future work
(\cref{out:oscillation}).

\begin{figure}[t]
  \centering
  \includegraphics[width=\textwidth]{\figdir/fig_zebrafish_artr_kinograph_decoders.png}
  \caption{\textbf{Recorded versus modelled ARTR across the three heading
  decoders.} Recorded ARTR neurons matched to their model counterparts and
  grouped by hemisphere (L block below the white divider), the model driven by
  the recorded $\omega$ and GCaMP-convolved. Columns: recorded $\Delta F/F$,
  then the model under the linear $(\cos\theta,\sin\theta)$, $64$-bin, and
  nonlinear MLP decoders.
  \textbf{(a)}~Recorded angular drive $\omega(t)$, a square-wave rotation
  command, shared across columns.
  \textbf{(b)}~Per-row $z$-scored kinographs; $r_{LR}$ is the L/R
  population-mean correlation.
  \textbf{(c, d)}~Per-hemisphere band decomposition of the population mean ---
  ARTR$_L$ (\textbf{c}, red) and ARTR$_R$ (\textbf{d}, blue) --- each an offset
  stack of the full trace (sum), its slow band (period $>16$\,s, the
  turn-block) and its fast band ($<16$\,s, the bump carrier), so that
  sum $=$ slow $+$ fast. The slow- and fast-band L/R correlations are annotated
  above each ARTR$_L$ panel: the recorded fast band is incoherent
  ($r_{LR}=+0.03$) whereas the modelled fast band is L/R phase-locked.}
  \label{fig:zhd_artr_kino}
\end{figure}

\begin{figure}[t]
  \centering
  \includegraphics[width=0.85\textwidth]{\figdir/fig_zebrafish_artr_oscillation_decoders.png}
  \caption{\textbf{ARTR oscillation diagnostic across the three heading
  decoders} (sign-constrained $917$-cell model, recorded $\omega$; the linear
  $(\cos\theta,\sin\theta)$, $64$-bin and MLP decoders overlaid).
  \textbf{(a)}~Dominant ARTR oscillation period versus constant $\omega$,
  against the bump-rotation period $360/|\omega|$ (dashed).
  \textbf{(b)}~ARTR oscillation amplitude (std of the population mean) versus
  constant $\omega$. The period tracks the bump and the amplitude stays
  roughly constant for all three decoders --- the inherited-carrier signature,
  independent of the readout.}
  \label{fig:zhd_artr_oscillation}
\end{figure}

\begin{figure}[t]
  \centering
  \includegraphics[width=\textwidth]{\figdir/fig_zebrafish_artr_dark.png}
  \caption{\textbf{ARTR at the rotation$\to$dark transition.} The continuous
  ZAPBench \emph{open loop} $\to$ \emph{rotation} $\to$ \emph{dark} blocks; the
  model is driven with the recorded rotation $\omega$ followed by zero input
  (dark block, shaded). \textbf{(a)}~Angular drive $\omega(t)$: zero in open
  loop, the $\pm30^\circ/$s square wave in rotation, zero in dark.
  \textbf{(b)}~Recorded ARTR L/R population mean (ARTR$_L$ red, ARTR$_R$ blue),
  normalised. \textbf{(c)}~Model ARTR L/R population mean, normalised.
  \textbf{(d, e)}~Recorded and modelled ARTR kinographs (per-row $z$-scored, L
  block below the white divider; dashed lines mark the rotation and dark
  onsets).}
  \label{fig:zhd_artr_dark}
\end{figure}


\section{Discussion}
\label{sec:discussion}

Taken together, the four results establish a single claim: a circuit whose
recurrent wiring is \emph{fixed} to the measured larval connectome (only
the edge magnitudes, the afferent gate, and the readout are free) can be
trained to perform the navigational computations the inferior--peduncular
region is thought to support, and the resulting model is informative about
\emph{how} the wiring does so.

First, the sign-locked circuit learns a near-unit-gain heading integrator
on the nominal connectome (\cref{sec:results_evolution}): a single
travelling activity peak tracks heading across held-out swims and
constant-drive sweeps. That integrator is a property of the recurrent
substrate, not of the observer reading it out: a localist ring, a
linear, and a nonlinear readout all recover the same dynamics, the only
caveat being that an over-flexible readout must be paced so the integrator
forms first (\cref{sec:results_decoder}). The same circuit, given two
anatomically distinct angular afferents, maintains a running
sensory-vs-proprioceptive \emph{difference} rather than echoing either
stream, a connectome-locked forward-model error signal
(\cref{sec:results_mismatch}).

Second, the circuit integrates heading together with translation, but
per-frame accuracy alone cannot decide \emph{which} translation the larval
brain performs: leaky and cumulative integrators, and 1-D distance versus
2-D path, all fit the swim statistics about equally well, so the choice
must be made on biological grounds rather than on fit quality
(\cref{sec:results_variants}). The biology points two ways. Leaky versus
cumulative is not a binary: a perfectly cumulative integrator demands a
marginally stable line attractor with exactly unit feedback, which is not
physically realisable \citep{seung1996,khona2022}, and the persistence time
constant of the goldfish and zebrafish oculomotor integrators is itself
detuned and retuned by feedback \citep{major2004,goncalves2014}; our
cumulative runs already drift toward a leaky solution, so the realistic
configuration is asymmetric: a near-lossless, topologically stabilised
heading ring \citep{kim2017,chaudhuri2019} alongside finite-time-constant,
leaky translation. The 2-D-versus-distance axis stays open: only the 2-D
target couples translation to the internally maintained heading, the
defining feature of true path integration, demonstrated by vector summation
in the fly central complex \citep{lyu2022}; the model's 2-D success
therefore states what the wiring \emph{can compute}, not that this circuit
holds an allocentric place code, which is established in the larval
telencephalon \citep{yang2024} but not in the hindbrain circuit we model.
Third, the recurrent code that supports this
is distributed and saturated: heading is encoded redundantly down to single
cells, the 2-D position only collectively (no neuron carries it, yet the
population decodes it near the information ceiling), and the split between
the two is set by the recurrent computation, not read off the cleanly
separated afferents, with pooled information growing only logarithmically in
neuron number (\cref{sec:results_four_classes}). Fourth, solving the task is
not the same as reproducing the brain: a connectome plus a task loss leaves
a genuine activity degeneracy that a calcium-observation loss only partially
closes, with the fish's own unmodelled behaviour setting the ceiling
(\cref{sec:results_task_vs_activity}).

These four results are one finding seen four times. The connectome fixes
\emph{that} the circuit integrates self-motion --- the topology and signs of
the recurrent pool make a heading integrator the natural solution --- but it
underdetermines almost everything else about the implementation. It
underdetermines the \emph{weights}: many sign-consistent edge-magnitude
configurations fit the same task equally well (\cref{sec:results_decoder}), and
the calcium-observation loss narrows but does not collapse this fold
(\cref{sec:results_obs_baseline}). It underdetermines the \emph{integrator
form}: leaky versus cumulative, and $1$-D distance versus $2$-D path, are
indistinguishable on per-frame fit and must be decided on biological grounds
(\cref{sec:results_variants}). It underdetermines the \emph{readout}: a
localist ring, a linear, and a nonlinear decoder all recover the same
underlying dynamics (\cref{sec:results_decoder}). And it underdetermines the
\emph{activity}: a circuit that solves the task perfectly can still produce
calcium quite unlike the recording (\cref{sec:results_task_vs_activity}). These
are not four separate caveats but four faces of a single underdetermination ---
the connectome constrains the computation far more tightly than it constrains
the weights, the dynamical form, the code that reads it, or the activity that
realises it. Naming and measuring that gap, rather than any single
task-solving result, is what this paper contributes.

\section{Outlook: which circuit properties are necessary, sufficient,
and biologically meaningful?}
\label{sec:outlook_future}

Three findings already in this paper are, we believe, its most
consequential, and each is a falsifiable structural prediction rather than a
reproduction of known physiology:

\begin{itemize}
\item \textbf{The afferents are separated; the recurrent code is mixed.}
  \Citet{feierstein2023} see distinct rotation and translation populations in
  the imaged hindbrain. The model places that separation at the \emph{input}:
  the angular and translational drives enter through anatomically distinct,
  sign-locked ports, while the recurrent integrator downstream mixes them onto
  a single low-dimensional manifold whose neurons carry both variables with
  only a graded bias (\cref{sec:results_four_classes}). The clean peripheral
  split is therefore an afferent property, not a property of the computation it
  feeds. This yields a direct prediction: imaging deeper into the recurrent
  pool should reveal mixed heading/translation selectivity, not the clean
  rotation/translation segregation reported at the periphery. The result
  reframes the imaging account rather than merely reproducing it.
\item \textbf{The two-dimensional translation code has no single-cell
  signature yet decodes near ceiling.} No neuron carries the planar position
  (per-neuron translation information $\approx\!0$), yet the population recovers
  it at $\approx\!3.4$ bits, on a low-variance subspace nearly orthogonal to
  the dominant heading ring and invisible to variance-ranked dimensionality
  reduction (\cref{fig:zhd_distributed_code}). This is a
  concrete, surprising claim about how a vertebrate circuit can hide a position
  code in plain sight, and it comes with its own test: the same subspace
  geometry should be recoverable from the recorded recurrent activity, not only
  from the trained model.
\item \textbf{The opponent code is reproduced; the intrinsic oscillator is
  not.} The connectome-locked circuit captures the ARTR's left/right opponent
  organisation but reproduces its fast alternation only while driven, going
  near-flat in darkness where the recording keeps oscillating ($|L\!-\!R|$ std
  $0.72$ recorded vs.\ $0.14$ model; \cref{fig:zhd_artr_dark}). The dissociation
  names the missing biology --- recurrent left$\leftrightarrow$right inhibition
  plus a slow adaptation variable --- and is a negative result that specifies a
  mechanism, which constrains the circuit more sharply than a positive
  capability claim would.
\end{itemize}

These three points motivate the directions below. The results otherwise leave
open the harder question of \emph{what about this particular network} is
actually doing the work. Nine directions follow directly from the experiments
already in this paper and together would convert the present existence proof
into a structural account of the computation.

\paragraph{Is the connectome necessary, and at what E/I ratio?}\label{out:necessary}\label{out:ei}
Re-training the same task family on random networks matched for size,
density, weight scale and E/I ratio (and on systematically rewired
versions of the measured connectome) isolates how much of the observed
performance comes from the wiring as opposed to the network's capacity. The
expected gradient (intact $\to$ degree-preserved rewiring $\to$
Erd\H{o}s--R\'enyi at matched density) speaks to whether the larval brain's
connectivity is computationally efficient or merely sufficient. The same
nulls also probe the E/I dimension: the recurrent pool is predominantly
inhibitory, and the v1 vs v2 connectome variants flip the \emph{IPN12} sign
while leaving every other edge fixed. Sweeping a controlled inhibitory
fraction at constant density and re-fitting the joint swim-integration task
would identify the window inside which both heading and forward distance
remain decodable, and how close the measured ratio sits to its centre. We
build these size-, density- and weight-scale-matched random null
connectomes (an inhibitory-fraction sweep together with ten replicates
matched to the measured E/I ratio) in \cref{app:er_null}
(\cref{fig:zhd_er_connectomes}).

\paragraph{What structure carries the computation --- how sparse, and
how divided?}\label{out:sparsity}\label{out:mixed} Two complementary
cuts probe how the trained recurrent pool actually does the work. An
edge-pruning sweep against the trained both / position-2D models, scoring
the task error as edges are removed in order of magnitude, gives the
sparsity floor below which integration fails, distinguishes a smoothly
degrading regime (graceful) from a brittle one (a cliff at a critical
sparsity), and identifies the small core of edges that carries the
computation. A finer per-neuron mutual-information analysis on
$\hat\theta$, $\hat d$, and their interaction (refining the coarse
angular/displacement partition of \cref{sec:results_four_classes})
would then say whether that core resolves into disjoint heading and
distance subpopulations or relies on neurons that jointly encode both,
and whether mixed selectivity emerges only when the task complexity
demands it.

\paragraph{How much does calcium supervision reduce solution
degeneracy?}\label{out:degeneracy} The observation loss already shifts each model toward a
small neighbourhood of recordings-consistent solutions
(\cref{sec:results_obs_sweep}). Training $K$ independent seeds with
$\lambda_{\mathrm{obs}} = 0$ and again with $\lambda_{\mathrm{obs}}
> 0$ and comparing the spread of $W^{\mathrm{rec}}$, $W^{\mathrm{in}}$,
$W^{\mathrm{out}}$, and the resulting per-cell tuning would quantify
how much identifiability the observation loss adds beyond the
connectome support and the task loss alone.

\paragraph{Exteroceptive vs proprioceptive translation.}\label{out:proprio} The
proprioceptive variant (\cref{method:proprioception}) routes
$v_{\mathrm{fwd}}$ through motor-efferent cells in parallel with the
exteroceptive (pt-IPN1) pathway. Ablating each pathway in turn at
test time (on the same trained circuit) decomposes the
translation drive into the two physiologically distinct sources and
predicts which calcium signatures should disappear under selective
silencing of either limb, providing a directly testable bridge to
in-vivo perturbation experiments.

\paragraph{What \emph{is} the distributed translation code?}\label{out:manifold}
\Cref{sec:results_four_classes} now answers the geometric form of this
question: the place estimate lives on a low-dimensional manifold (participation
ratio $3$--$5$) as a low-variance translation subspace nearly orthogonal to a
dominant heading ring, rank-2-readable by a task-tuned but not a
variance-aligned readout, with the heading ring reused and the translation
subspace rotated as the task moves from $1$-D distance to a $2$-D path; the
same reduced-rank analysis \citet{feierstein2023} applied to the imaged
hindbrain recovers separate rotation and translation modes at the afferents
that the recurrent integrator then mixes. What remains open is the
\emph{dynamical} side of this geometry: how trajectories move on the manifold
during a swim bout, whether the translation subspace is a true line / plane
attractor or a transient, and (the decisive test) whether the same
subspace geometry is recoverable from the \emph{recorded} dIPN activity rather
than only from the trained model, which would tie the account to the biology
directly.

\paragraph{Do latent translation variables emerge without translation
supervision?}\label{out:latent_translation}
The geometry analysis suggests that translation occupies a population subspace
distinct from the dominant heading ring, but it remains unclear whether this
representation is induced by translation supervision or already present in the
connectome-constrained dynamics. A direct control is to train the circuit on
heading alone while retaining the forward-velocity input and then decode
distance and position from the recurrent state. Preliminary experiments
indicate that a heading-only model already carries partially decodable distance
and 2-D position information despite never being supervised on either quantity.
However, current integral-versus-leaky-filter tests favour a short-timescale
trace rather than a stable accumulated estimate, suggesting that the latent
signal may reflect a recurrent memory of forward velocity rather than a true
path integrator. Resolving this distinction requires fully converged
heading-only models, replication across random seeds, and comparison against
matched Erd\H{o}s--R\'enyi null connectomes. Together, these controls would
distinguish translation signals inherited from the inputs, translation
representations induced by supervision, and latent computations that emerge
naturally from the measured recurrent wiring.

A further sharpening tracks this signal as training proceeds. Because any
recurrent network passively low-pass-filters its forward-velocity input, a
leaky-filtered distance is decodable from the recurrent state \emph{by
construction} --- a matched random null already yields it at initialisation ---
so a positive explained variance for the leaky trace is uninformative on its
own. The discriminating quantity is instead the \emph{gap} between the explained
variance of the true cumulative integral and that of the leaky filter,
$R^2_{\mathrm{int}}-R^2_{\mathrm{leaky}}$, followed as a function of training
epoch. The signature of a genuinely emergent integrator would be that this gap
\emph{grows} over the course of heading training in the measured connectome but
not in the matched random control: the passive velocity echo present at
initialisation giving way to a stable accumulated estimate that the specific
measured wiring, and not generic recurrence, is able to build.

\paragraph{How and why are the fast oscillations generated --- and do they
drive heading?}\label{out:oscillation} The recorded ARTR carries a fast hemispheric alternation
that the connectome-locked model reproduces only while driven and loses in
darkness (\cref{fig:zhd_artr_oscillation,fig:zhd_artr_dark}): in the model
the oscillation is \emph{inherited}, not intrinsic. Whether the biological
rhythm is generated locally (by recurrent left$\leftrightarrow$right
inhibition plus a slow adaptation variable) or inherited from upstream,
and, crucially, whether it \emph{drives} the heading estimate or merely
rides on it, is unresolved. Adding the two missing ingredients to the
connectome-locked circuit and asking whether a self-sustained oscillation
then emerges, and whether suppressing it degrades integration, would test
both its origin and its causal role.

\paragraph{Can the recording be recovered by injecting a learned drive?}\label{out:inr}
The calcium reconstruction is bounded from above by signals the model never
sees, chiefly the fish's own behaviour (\cref{sec:results_unmodeled}).
Rather than hand-craft those inputs, one could \emph{learn} the missing
drive as a continuous implicit neural representation (an INR mapping time,
or behavioural state, to an injected proprioceptive / efference signal)
trained jointly with the observation loss so that the connectome-locked
circuit reproduces the recorded activity. If a low-dimensional learned
drive closes the remaining gap, it would both name the behavioural variable
the circuit is tracking and give a generative model of the recording; if it
does not, the residual marks structure the connectome-locked dynamics
genuinely cannot produce.

\paragraph{Why is heading easily decoded in the model but not in the
recording?}\label{out:decodability}
One of the most striking discrepancies in this study is that heading is almost
perfectly decodable from the connectome-constrained models yet remains difficult
to recover from the recorded ZAPBench calcium activity. At least four
explanations are possible. First, the observable itself may destroy much of the
signal: calcium kinetics, low imaging rates, and partial observation of the
circuit substantially degrade heading decoding even in the model. Second,
heading may be represented dynamically rather than instantaneously, so that it
is recoverable from population trajectories or temporal windows but not from
single calcium frames. Third, the stimulus-defined heading used in the model may
not be the variable represented most faithfully by the biological circuit, which
could instead track internal estimates of self-motion, motor state, uncertainty,
or other latent quantities. Fourth, the model may be too clean: training
explicitly optimises heading, whereas the biological population simultaneously
represents multiple sensory, motor, and internal-state variables, diluting any
single decodable signal. Distinguishing between these possibilities requires
moving beyond frame-by-frame decoding and comparing alternative behavioural and
latent variables against the recorded activity. More generally, this discrepancy
highlights the central inverse-problem question of the paper: a computation may
be obvious in a mechanistic model yet remain difficult to identify from the
experimentally accessible observable.

\appendix
\section{Supplementary figures: training dashboards}
\label{app:dashboards}
The connectome-locked training dashboards for the model variants of \cref{sec:results}, in the same panel layout as \cref{fig:zhd_evolution_artr_pt1}.

\begin{figure}[p]
  \centering
  \includegraphics[width=0.9\textwidth]{\figdir/fig_evolution_zebrafish_hd_si_ipn_917_v1_selfmotion_rotation_kbins.png}
  \caption{\textbf{Heading integrator with a $64$-bin softmax decoder.}
  Same circuit and heading task as \cref{fig:zhd_evolution_artr_pt1},
  with the linear $(\cos\theta,\sin\theta)$ readout replaced by a
  $64$-bin heading-bin readout ($n_{\mathrm{out}}=64$). Panels
  \textbf{(a)}--\textbf{(e)} follow \cref{fig:zhd_evolution_artr_pt1}
  (GT and learned $\hat W^{\mathrm{rec}}$, per-edge weight scatter,
  all-neuron and dIPN kinographs). \textbf{(f)}~Constant-$\omega$
  rollout: $\omega(t)$ and the bin-decoded heading vs.\ truth.
  \textbf{(g)}~Held-out swim trial, same overlay. \textbf{(h)}~The
  $64$-bin softmax readout over the rollout (bin index vs.\ time; the
  diagonal band is the travelling heading estimate, GT bin overlaid).
  \textbf{(i)}~Five held-out trials and \textbf{(j)}~five constant-$\omega$
  sweeps, bin-decoded heading vs.\ truth.}
  \label{fig:zhd_evolution_kbins}
\end{figure}

\begin{figure}[p]
  \centering
  \includegraphics[width=0.9\textwidth]{\figdir/fig_evolution_zebrafish_hd_si_ipn_917_v1_propriocep_mismatch.png}
  \caption{\textbf{Trained proprioceptive-gain mismatch integrator.}
  Sign-locked $917$-cell RNN on the proprioception circuit
  (\cref{fig:zhd_3d}\textbf{d}), driven by the $5$-channel proprioception
  input and supervised on $(\cos\theta,\sin\theta,m)$ with $m$ the
  integrated mismatch (\cref{method:mismatch}). Panels
  \textbf{(a)}--\textbf{(e)} and \textbf{(h)} follow
  \cref{fig:zhd_evolution_artr_pt1} (GT and learned
  $\hat W^{\mathrm{rec}}$, per-edge weight scatter, all-neuron and dIPN
  kinographs, integration-gain scatter).
  \textbf{(f)}~Constant-$\omega$ deterministic rollout: $\omega(t)$ and
  decoded vs.\ true heading.
  \textbf{(g)}~Mismatch rollout --- top: the sensory $\omega$ and the
  proprioceptive $\omega_{\mathrm{proprio}}=g\,\omega$; middle: the two
  diverging integral paths $\int\omega$ and $\int\omega_{\mathrm{proprio}}$;
  bottom: decoded (black) vs.\ true (green) integrated mismatch $m$.
  \textbf{(i)}~Five held-out swim trials, decoded vs.\ true heading.
  \textbf{(j)}~The same trials, decoded vs.\ true integrated mismatch
  $m$ --- the companion of panel \textbf{g} over multiple trials.
  \textbf{(k)}~Constant-$\omega$ deterministic heading sweeps.}
  \label{fig:zhd_evolution_mismatch}
\end{figure}

\begin{figure}[p]
  \centering
  \includegraphics[width=0.9\textwidth]{\figdir/fig_evolution_zebrafish_hd_si_ipn_917_v1_selfmotion_both.png}
  \caption{\textbf{Trained heading $+$ distance integrator,
  cumulative target.} Same panel layout as
  \cref{fig:zhd_evolution_artr_pt1}; panel (g) overlays the GT and
  decoded forward-distance trace below the heading components.}
  \label{fig:zhd_evolution_both}
\end{figure}

\begin{figure}[p]
  \centering
  \includegraphics[width=0.9\textwidth]{\figdir/fig_evolution_zebrafish_hd_si_ipn_917_v1_selfmotion_both_leaky.png}
  \caption{\textbf{Trained heading $+$ distance integrator,
  leaky target} ($\tau_d=0.5\,$s). Same panel layout as
  \cref{fig:zhd_evolution_both}.}
  \label{fig:zhd_evolution_both_leaky}
\end{figure}

\begin{figure}[p]
  \centering
  \includegraphics[width=0.9\textwidth]{\figdir/fig_evolution_zebrafish_hd_si_ipn_917_v1_position_2d.png}
  \caption{\textbf{Trained 2-D path integrator, cumulative target.}
  Panels (a)--(e), (h)--(i) follow \cref{fig:zhd_evolution_artr_pt1};
  panels (f)--(g) show the decoded path $(x,y)$ in spatial coordinates
  on the constant-$\omega$ rollout and on a held-out swim trial.}
  \label{fig:zhd_evolution_position2d}
\end{figure}

\begin{figure}[p]
  \centering
  \includegraphics[width=0.9\textwidth]{\figdir/fig_gnn_dashboard_zebrafish_hd_si_gnn_ipn_917_v1_selfmotion_rotation.png}
  \caption{\textbf{Message-passing GNN on the heading-only task.}
  Companion of \cref{fig:zhd_evolution_artr_pt1}. Same circuit, gate,
  and heading head as the sign-locked RNN, with the dense recurrent
  operator replaced by per-edge and per-neuron MLPs over a learnable
  per-neuron embedding.
  \textbf{(a)}~GT connectome $W^{\mathrm{con}}$, neuron-type blocks on
  both axes. \textbf{(b)}~Learned per-edge gain on the same scale.
  \textbf{(c)}~Scatter of the learned per-edge gain against
  $W^{\mathrm{con}}$. \textbf{(d)}~All $917$ neurons over a
  constant-$\omega = +60^\circ/\mathrm{s}$ rollout, sorted by cell type.
  \textbf{(e)}~dIPN sub-population over the same rollout, rows ordered by
  preferred heading. \textbf{(f)}~Constant-$\omega$ rollout: $\omega(t)$
  and decoded (black) vs.\ true (green) heading. \textbf{(g)}~A held-out
  swim test trial, same overlay. \textbf{(h)}~Integration-gain scatter:
  measured slope of the decoded heading vs.\ true $\omega$.
  \textbf{(i)}~The per-neuron embedding, coloured by cell type.
  \textbf{(j)}~The learned per-edge signalling function (mean $\pm$ SD
  per cell type) over the subthreshold-state range. \textbf{(k)}~The
  learned node update at zero recurrent input (mean $\pm$ SD per cell
  type). Panels (i)--(k) are the GNN-specific row absent from the RNN
  figure.}
  \label{fig:zhd_gnn_dashboard}
\end{figure}

\begin{figure}[H]
  \centering
  \includegraphics[width=0.9\textwidth]{\figdir/fig_zebrafish_mi_partition_gnn.png}
  \caption{\textbf{Pooled mutual information for the message-passing GNN
  models} --- the companion of \cref{fig:zhd_mi_partition} for the GNN
  variant (\cref{method:GNN}; per-edge $g_\phi$ and per-neuron $f_\theta$
  MLPs on the same connectome, gate, and heads). Same panels, axes, and
  dot-size encoding as \cref{fig:zhd_mi_partition}:
  \textbf{(a)}~heading\,$+$\,distance (cumulative), \textbf{(b)}~leaky,
  \textbf{(c)}~2-D path (cumulative), \textbf{(d)}~leaky.}
  \label{fig:zhd_mi_partition_gnn}
\end{figure}

\begin{figure}[H]
  \centering
  \includegraphics[width=0.9\textwidth]{\figdir/fig_zebrafish_mi_grouped.png}
  \caption{\textbf{Logarithmic redundancy law of the pooled code.} For each
  of the four joint-target models, the cross-validated group decoder
  information (in bits, vertical axis) is plotted against the summed
  single-neuron information of the same group (horizontal axis, logarithmic),
  with one point per cell type and per anatomical family and a cross for the
  whole recurrent pool. The group information rises only logarithmically with
  the pooled single-cell information, a shallow slope that is the quantitative
  statement of diminishing returns: each added neuron contributes little once
  the pool is large. The overlaid fit reports the per-decade slope and its
  goodness of fit; heading and translation are shown separately.}
  \label{fig:zhd_mi_grouped}
\end{figure}

\begin{figure}[p]
  \centering
  \includegraphics[width=0.9\textwidth]{\figdir/fig_evolution_zebrafish_hd_si_ipn_917_v1_selfmotion_rotation_mlpdec.png}
  \caption{\textbf{Trained heading integrator with a nonlinear MLP decoder.}
  The same $917$-cell circuit and heading task as
  \cref{fig:zhd_evolution_artr_pt1}, with the linear
  $(\cos\theta,\sin\theta)$ readout replaced by a $2$-layer, width-$64$
  $\tanh$ MLP, and the recurrent core paced to lead the decoder
  (\cref{sec:results_decoder}). Panels follow
  \cref{fig:zhd_evolution_artr_pt1}: \textbf{(a--c)} ground-truth and
  learned recurrent weights and the per-edge scatter; \textbf{(d, e)}
  all-neuron and dIPN kinographs over a constant-$\omega$ rollout;
  \textbf{(f, g)} constant-$\omega$ and held-out swim rollouts (drive on
  top, decoded vs.\ true heading below); \textbf{(h)} integration-gain
  scatter; \textbf{(i)} five held-out swim trials and \textbf{(j)} five
  constant-$\omega$ sweeps, decoded (black) vs.\ true (green) heading with
  per-panel circular RMSE and Pearson $r$.}
  \label{fig:zhd_evolution_mlpdec}
\end{figure}


\section{Randomized null-connectome controls}
\label{app:er_null}

To make the necessity question of \cref{sec:outlook_future} concrete we
construct random null connectomes against which the same heading task can be
re-fitted, in \emph{three} families that isolate progressively finer aspects of
the wiring. All of them are matched to the measured $917$-cell circuit on the
number of nodes ($N=917$), the edge count (and hence density: $30{,}851$ edges,
density $0.037$) and the weight scale --- each random magnitude drawn
independently from the empirical distribution of measured synaptic weights ---
and all inherit the \emph{node identities} unchanged (cell-type assignment,
afferent gate targets, recurrent ring ordering), so every null remains a valid
drop-in substrate for the same trained model. The families differ in how much of
the measured connectivity they preserve.

\paragraph{Full-network rewiring, inhibitory-fraction sweep.} The first family
draws a brand-new random support uniformly over the whole adjacency matrix,
destroying the block structure entirely. Following Dale's law each cell keeps a
single output polarity; the afferent populations stay excitatory while a
controlled fraction $p_{\mathrm{inh}}$ of the recurrent pool is made inhibitory.
The four settings $p_{\mathrm{inh}}\in\{0.4,0.6,0.8,1.0\}$ span overall
inhibitory edge fractions of $30.7\%$, $46.0\%$, $61.4\%$ and $76.8\%$,
bracketing the measured $66.3\%$, and probe the inhibitory-fraction window of
\cref{sec:outlook_future}.

\paragraph{Full-network rewiring, E/I matched.} The second family is the same
full-network rewiring but with its overall inhibitory edge fraction matched
\emph{exactly} to the measured $66.3\%$; ten independent draws (different random
support and polarity assignment) give the null-model variance against which the
intact circuit is read. This is the necessity control proper: it asks whether
\emph{any} random network of the measured size, density, weight scale and
excitation/inhibition balance reproduces the computation, or whether the
specific measured wiring is required.

\paragraph{Block-preserving shuffle.} The third family is the intermediate null.
Instead of a fresh random support it keeps the measured block-to-block
connectivity exactly --- every functional block-pair retains its edge count and
its magnitude distribution --- and randomises only \emph{which} cells within
those blocks are wired. The coarse block architecture (which functional
populations connect, and how densely) is therefore identical to the measured
circuit while the microscopic wiring is scrambled; ten independent draws again
give the null band. Where the full-network families ask whether the wiring
matters \emph{at all}, the block-preserving family separates the contribution of
the coarse block architecture from that of the fine within-block wiring: a
computation that survives the block-preserving shuffle but fails the
full-network nulls is carried by the block-level organisation, whereas one that
fails even the shuffle depends on the precise cell-to-cell connectivity.

\Cref{fig:zhd_er_connectomes} renders the measured connectome alongside one
representative of each family in the same $z$-scored, partition-sorted
signed-matrix format as \cref{fig:zhd_connectome_summary}: the measured block
structure is visibly retained by the block-preserving shuffle (block-density
correlation $\approx1$) but erased by both full-network rewirings ($\approx0$),
with the block-preserving and E/I-matched families at the measured $66.3\%$
inhibitory fraction and the inhibitory-fraction sweep spanning it by design.
Re-fitting the heading task on each family and
comparing the recovered integration accuracy against the intact circuit is the
experiment of \cref{sec:outlook_future}; the matched random substrates and the
training configuration are prepared here.

\begin{figure}[p]
  \centering
  \includegraphics[width=\textwidth]{\figdir/fig_er_connectomes.png}
  \caption{\textbf{Three null-connectome families vs.\ the measured circuit.}
  \textbf{(a--d)}~Same $z$-scored, partition-sorted signed-matrix rendering as
  \cref{fig:zhd_connectome_summary} (red excitatory, blue inhibitory; partition
  blocks on both axes), for the measured circuit and one representative of each
  null family. \textbf{(a)}~The measured $917$-cell connectome, block-structured
  by the functional partition. \textbf{(b)}~A full-network random rewiring
  (inhibitory-fraction sweep): the block structure is destroyed.
  \textbf{(c)}~A full-network rewiring matched to the measured
  excitation/inhibition balance: still structureless, with matched colour
  balance. \textbf{(d)}~A block-preserving shuffle: the block-to-block
  connectivity is kept while the fine within-block wiring is randomised. All
  nulls match the measured size, density and weight scale and inherit the node
  identities, so the excitatory afferent columns (right) are unchanged
  throughout. Quantitatively, the block-preserving shuffle reproduces the
  measured block-density map almost exactly (block-density correlation
  $\approx1$) at the measured $66.3\%$ inhibitory fraction, whereas both
  full-network rewirings erase the block structure (correlation $\approx0$).}
  \label{fig:zhd_er_connectomes}
\end{figure}


\section{Latent translation probe}\label{app:latent_translation}
This appendix collects the control experiment proposed in
\cref{out:latent_translation}: training the circuit on heading alone while
keeping the forward-velocity input, then asking whether the recurrent state
nonetheless carries a translation estimate, and if so whether that estimate is a
genuine integral or a passive low-pass echo of the velocity drive, and whether
it depends on the measured wiring. \Cref{fig:zhd_vfwd_distance} compares the
measured circuit against a randomized-connectome control matched on size,
density, weight scale and excitation/inhibition balance, with only the recurrent
topology shuffled. For each one a linear decoder is fit from the recurrent state
to the forward distance on held-in trials and applied to held-out trials,
against two references: the true cumulative integral (which a genuine integrator
tracks) and a leaky-filtered velocity (which a passive echo saturates onto). A
state that follows the cumulative reference, on the measured circuit but not the
randomized one, would identify the latent translation estimate as a genuine
path-integration variable built by the specific connectome.

\begin{figure}[H]
  \centering
  \includegraphics[width=\textwidth]{\figdir/fig_zebrafish_vfwd_distance_dashboard.png}
  \caption{\textbf{Latent forward-distance decode, measured connectome vs.\
  randomized control.} Rows are held-out swim trials. The four columns are the
  two heading-only models crossed with the two ground-truth references the
  decoder is fit to: \textbf{(a)}~measured circuit, cumulative integral;
  \textbf{(b)}~measured circuit, leaky-filtered velocity; \textbf{(c)}~randomized
  connectome, cumulative integral; \textbf{(d)}~randomized connectome, leaky
  velocity. Within a column every panel reads out the \emph{same} recurrent
  state; columns differ only in the reference the linear decoder is regressed on,
  not in the model (which is never supervised on distance). In each panel the
  green trace is the ground-truth reference and the black trace is its decode
  from the recurrent state, with the per-trial explained variance annotated. A
  decode that tracks the cumulative reference (column a / c) indicates a genuine
  integral; one that only tracks the leaky reference (column b / d) indicates a
  passive low-pass echo. The measured-vs-randomized contrast tests whether any
  property recovered here is specific to the measured wiring.}
  \label{fig:zhd_vfwd_distance}
\end{figure}


\section{ZAPBench training data}\label{method:zapbench}
The recorded activity is the publicly released ZAPBench whole-brain
light-sheet calcium recording of a larval zebrafish, registered to the
same \emph{fish2} EM volume as the connectome (\cref{sec:anatomy}): the
$\Delta F/F$ trace of every segmented cell is taken from the ZAPBench
release and each connectome \emph{bodyId} is matched to its functional
ROI through the EM\,$\leftrightarrow$\,functional registration. Of the
$839$ neurons in the IPN12 circuit, $481$ carry a matched recording:
$300$ ring-pool cells (\emph{IPNd}/\emph{IPNds}/\emph{IPN12}) and $181$
afferents (\emph{RIPN}, \emph{pt-IPN}). \Cref{fig:zhd_allblocks} lays out
the whole $\sim\!2\,\mathrm{h}$ recording: the nine task blocks, the fish's
own swim behaviour, the recorded activity of the three velocity-gate
afferent classes (ARTR, pt-IPN1, motor\_efferent), and the ring pool,
the same $481$ matched neurons used throughout. The observation loss is
taken over all $481$ matched cells; the real-vs-model kinographs display only
the $300$ ring-pool subset. (These matched-cell counts are for the
$839$-cell \emph{IPN12 precursor} on which the calcium observation-loss
experiments of \cref{sec:results_obs_sweep} were run, the circuit
\emph{before} the IPN-core families were added to form the nominal
$917$-cell circuit of \cref{sec:anatomy}. The task-only calcium baseline
of \cref{fig:zhd_obs_baseline} has since been re-run on the full
$917$-cell circuit, where $279$ ring-pool cells carry a matched
recording; the task-integration results of \cref{sec:results} all use
the $917$-cell circuit.) We supervise on the ``Rotations'' block, in which a
$45^\circ/\mathrm{s}$ drifting grating rotates the animal's virtual
heading: the true heading is read from the recorded stimulus channel and
differentiated to the angular-velocity drive $s(t)$, so the \emph{model}
is driven by the \emph{same} stimulus that produced the recording. The
block spans ${\sim}600\,\mathrm{s}$ ($656$ imaging frames at the
$1.09\,\mathrm{Hz}$, $0.915\,\mathrm{s}$ imaging rate). Because the model
runs at $\Delta t = 0.01\,\mathrm{s}$, the recorded $\Delta F/F$ is
linearly interpolated up to the model grid (edge-padded), and each neuron
is $z$-scored \emph{globally} over the whole block (not per window),
so the target keeps a single, trial-independent scale. The block is then
sliced into $1000$-frame ($10\,\mathrm{s}$) trials drawn from disjoint
train/test time spans, in the same format as the synthetic
swim-integration trials: channel~0 is the recorded drive $s(t)$ and the
first frame carries the $(\cos\theta_0,\sin\theta_0)$ heading cue, with no
warm-up. During training each task minibatch is augmented with $B_c$ of
these calcium trials in a single forward pass; the heading MSE of
\cref{eq:zhd_loss} is taken over the combined batch (the rotation trials
also carry a true heading), while $\mathcal{L}_{\mathrm{obs}}$ is
evaluated on the calcium trials only.

\begin{figure}[p]
  \centering
  \includegraphics[width=\textwidth]{\figdir/fig_zebrafish_all_blocks.png}
  \caption{\textbf{The full ZAPBench session: nine conditions, one circuit.}
  All locally-available drives and behaviour over the whole recording
  ($7870$ frames $\approx 2\,\mathrm{h}$; the same $481$ matched neurons used
  throughout), with the nine task blocks separated and labelled
  (\textsc{gain, dots, flash, taxis, turning, position, open loop, rotation,
  dark}). \textbf{Top:} the per-frame stimulus-condition code. \textbf{Next
  two:} the fish's own swim behaviour --- \emph{forward} (bilateral tail-EMG
  envelope) and \emph{turning} (L/R asymmetry) --- available continuously in
  every block. \textbf{Middle three:} the recorded $\Delta F/F$ of the three
  velocity-gate afferent classes, in the \cref{fig:zhd_3d} colour code ---
  \emph{ARTR} (blue, angular drive), \emph{pt-IPN1} (orange, exteroceptive
  forward) and \emph{motor\_efferent} (green, angular efference);
  matched cells only, $z$-scored per neuron. \textbf{Bottom:} the recorded
  ring-pool $\Delta F/F$ ($n{=}300$, rastermap order, $z$-scored; same row set
  as \cref{fig:zhd_obs_reconstruction,fig:zhd_unmodeled}). We have so far
  trained only on the Rotation slice.}
  \label{fig:zhd_allblocks}
\end{figure}

\section{Reconstructing the recorded calcium activity}
\label{app:calcium}

\subsection{Task ability does not imply realistic activity}
\label{sec:results_obs_baseline}

The results so far establish that the connectome-constrained network learns
to integrate heading accurately, and that the solution is robust to any
single-type lesion. Accuracy on the task, however, says nothing about
whether the network's \emph{internal activity} matches the brain. We test
this directly by driving the trained task-only model
(\emph{zebrafish\_hd\_si\_ipn\_917\_v1}, the nominal $917$-cell circuit,
trained on the rotation task \emph{only}) with the recorded ZAPBench
``Rotations'' stimulus (\cref{method:zapbench}), converting its
rolled-out activity to GCaMP7f calcium (\cref{eq:zhd_gcamp}), and
comparing it to the recorded $\Delta F/F$ of the $279$ matched ring-pool
neurons in the same rastermap row order (\cref{fig:zhd_obs_baseline}).

The model is an excellent head-direction predictor (its decoded heading
tracks the $45^\circ/\mathrm{s}$ rotation closely,
\cref{fig:zhd_obs_baseline}, bottom panel), yet its population activity is
plainly \emph{not} the recorded one: the model's learned kinograph is a
clean, stimulus-locked set of bands, whereas the recorded $\Delta F/F$
above it is irregular, bursty and drift-laden. This
is the degeneracy at the heart of the paper, now read out on real data:
\emph{many circuits realise the same head-direction computation, and hence
the same task ability, while producing different activity}. Task
supervision constrains only the two-dimensional readout
(\cref{eq:zhd_readout}); the $N$-dimensional state it leaves unobserved is
exactly where implementations diverge.

A natural objection is that the mismatch is merely cosmetic: perhaps the
recorded activity still encodes heading, just in a form our eye does not
read off the kinograph. We tested this directly by training a decoder
(a small neural network) to predict the instantaneous heading
$(\cos\theta,\sin\theta)$ from the recorded per-frame $\Delta F/F$ of the
ring-pool neurons (fit on the first $70\%$ of the rotation block,
evaluated on the held-out tail). The result was \emph{not} conclusive: the
decoded heading does not track the true heading
(\cref{fig:zhd_obs_control}; recorded-$\Delta F/F$ decode at circular
$r=0.14$), in sharp contrast to the model's own readout, which
recovers heading almost perfectly from its latent state
(\cref{fig:zhd_obs_baseline}, bottom panel). The implication is that \emph{the
instantaneous head-direction signal is hard to read, frame by frame, from
the observed calcium activity}.

How much of this failure is the \emph{observable} rather than the brain? We
settle the confound with a control: run the \emph{identical} decoder on the
trained model pushed through the \emph{same} observation pipeline: its
rolled-out activity convolved with the GCaMP kernel, sampled at the
$1.09\,\mathrm{Hz}$ imaging rate and restricted to the same $300$ cells
(\cref{fig:zhd_obs_control}). The model provably encodes heading: its latent
readout decodes the heading at circular $r = 0.97$. Through the degraded
observable the \emph{same} circuit decodes at only $r = 0.73$. The
observation pipeline alone (GCaMP low-pass of the fast $45^\circ/\mathrm{s}$
ring activity, $\tau_d = 1.2\,\mathrm{s}$, \cref{eq:zhd_gcamp}; $1.09\,\mathrm{Hz}$
aliasing of an $8\,\mathrm{s}$ revolution; only $300$ of $839$ cells seen)
removes a substantial fraction of the per-frame legibility. It does
\emph{not}, however, reach the real-data level ($r = 0.14$): a circuit that
encoded heading as cleanly as the model would still be decodable at
$r \approx 0.73$ under the recording's observation, so the observable is a
major but \emph{not} the sole cause. The residual gap reflects additional
real-world degradation: recording noise, behavioural variability of the
animal relative to the grating, and the idealised, stimulus-locked model
ring activity being crisper than the brain's. Either way heading is carried in the
population's \emph{temporal dynamics}, not in any single calcium frame.
This is why the observation loss matches the time-course of each neuron's
trace (\cref{eq:zhd_obs}, a per-neuron shape loss) rather than asking the
activity to decode heading frame-by-frame.

This motivates the calcium observation loss (\cref{eq:zhd_obs}): by adding
a term that pulls the model's predicted calcium toward the recorded
$\Delta F/F$, we aim to \emph{partially break} this degeneracy,
selecting, among the task-equivalent circuits, those whose activity is
also consistent with the measured brain. Runs that train with the
observation loss are reported in \cref{sec:results_obs_sweep};
\cref{fig:zhd_obs_baseline} is the no-observation-loss
($\lambda_{\mathrm{obs}} = 0$) baseline against which they are compared.

\begin{figure}[p]
  \centering
  \includegraphics[width=0.7\textwidth]{\figdir/fig_zebrafish_calcium_baseline.png}
  \caption{\textbf{Task-only model driven by the recorded Rotations
  stimulus and rendered as calcium.}
  The nominal $917$-cell model
  (\emph{zebrafish\_hd\_si\_ipn\_917\_v1}, heading task only, no
  observation loss) driven by the recorded ZAPBench Rotations stimulus
  and rendered as GCaMP7f calcium (\cref{eq:zhd_gcamp}), in the same
  reconstruction layout as \cref{fig:zhd_obs_reconstruction}. Top:
  recorded angular-velocity drive $\omega(t)$. For the bump-pool
  ($n=279$ matched ring-pool neurons, official ZAPBench rastermap row
  order) and the \emph{ARTR} afferent class ($n=33$, sorted by per-cell
  preferred-heading angle, \cref{fig:zhd_unmodeled}), two rows are
  shown: the recorded $\Delta F/F$ and the model's calcium under a
  continuous $600\,\mathrm{s}$ rollout, with per-panel $z$-MSE and SSIM
  inset; each row is the same recorded cell across the recorded and
  model panels. Bottom: true heading (green) and the model's
  readout-decoded heading (black) over the same rollout. This is the
  no-observation-loss baseline for the degeneracy the observation loss
  (\cref{eq:zhd_obs}; \cref{sec:results_obs_sweep}) targets; pt-IPN1 and
  motor\_efferent are omitted here (shown with the observation loss on in
  \cref{fig:zhd_obs_reconstruction}).}
  \label{fig:zhd_obs_baseline}
\end{figure}

\begin{figure}[t]
  \centering
  \includegraphics[width=\textwidth]{\figdir/fig_zebrafish_obs_control.png}
  \caption{\textbf{Observability control: the same per-frame decoder applied
  to the model's matched observable.} The task-only model is driven by the
  ZAPBench Rotations stimulus and its decoded heading read three ways over
  the block (circular correlation $r$ on the held-out tail).
  \textbf{Left:} true heading (green) vs.\ the model's latent readout (black);
  the MLP decode of the model's GCaMP-convolved,
  imaging-sampled, $300$-cell observable (red); and the MLP decode
  of the recorded $\Delta F/F$ (blue); the time axis is shown for
  the first $300\,\mathrm{s}$ for legibility (the $r$ values are computed over
  the full held-out tail). \textbf{Right:} the three held-out correlations
  as a bar chart.}
  \label{fig:zhd_obs_control}
\end{figure}

\subsection{Signals not yet accounted for: the fish's own behaviour}
\label{sec:results_unmodeled}

The task-only model's calcium is essentially uncorrelated with the
recording (\cref{sec:results_obs_baseline}; \cref{fig:zhd_obs_baseline}),
and the ceiling on what any model can recover here is set in part by
what the model is \emph{allowed to see}. The only time-varying input
to the network is the angular-velocity drive $s(t)$, which for the
Rotations block is an essentially periodic $\pm 45^\circ/\mathrm{s}$
square wave. A deterministic circuit driven by a periodic input can
only produce (after transients) a periodic output, so \emph{any}
non-periodic structure in the recorded $\Delta F/F$ is, by construction,
outside the model's reach. The recorded ring activity is
conspicuously non-periodic across the $600\,\mathrm{s}$ block even
though the stimulus repeats.

\Cref{fig:zhd_unmodeled} lays the recording's covariates side by side
over the rotation block to make the missing signals explicit. Above
the recorded $\Delta F/F$ we plot, from the raw ZAPBench ephys, the
fish's own swimming behaviour, the quantity the model never
receives. The two bottom traces are the left and right tail-muscle
EMG channels (the independent motor readout); the larval-zebrafish
swim decode summarises them as a \emph{forward} drive (bilateral
envelope) and a \emph{turning} signal (left/right asymmetry). These
behave very differently with respect to the stimulus. Turning is
largely stimulus-locked (it correlates with the drive at
$\rho \approx 0.76$, the expected optomotor following of the grating),
and so a stimulus-driven model can in principle capture it.
Forward drive, by contrast, is a slow, non-periodic latent: its
lag-one autocorrelation is $\approx 0.97$ and it is essentially
uncorrelated with the stimulus magnitude ($|\rho| \lesssim 0.05$). It
drifts on the tens-of-seconds-to-minutes timescale that dominates the
recurrent pool's irregular activity, yet nothing in $s(t)$ carries
it.\footnote{The forward/turning traces are derived from the two
tail-EMG electrodes by the standard ZAPBench swim-EMG pipeline
(windowed-variance envelope, swim-bout detection, per-bout weighting);
we verified they are genuinely those channels by recovering
\emph{forward} from the bilateral EMG sum ($\rho \approx 0.57$) and
\emph{turning} from the left$-$right difference ($\rho \approx 0.42$).
The per-frame RMS shown here is a coarse envelope of the same
channels.}

Three classes of signal are therefore present in the recording but
absent from the current model. (i) \emph{Self-generated motor drive}:
the forward/turning behaviour above, which is internally generated
and not a deterministic function of the visual stimulus; the forward
component in particular injects exactly the slow, aperiodic variation
the model cannot synthesise. (ii) \emph{Slow physiological and state
drift}: bleaching, arousal, and adaptation accumulate over the
block, which is why per-trial $z$-scoring is necessary in
\cref{eq:zhd_obs} to remove this component rather than fit it.
(iii) \emph{Unmodelled sensory channels}: the recording also
carries the raw visual-velocity and other stimulus channels at full
temporal resolution, of which the model receives only the distilled
heading drive. None of these are failures of identifiability: they
are inputs the network is simply not given. With this in mind we now
turn the observation loss on (\cref{sec:results_obs_sweep}) and see
how far it can close the gap.

\paragraph{Overall: do the results agree with the introduction?}
The introductory picture is a ring-attractor head-direction circuit:
hundreds of GABAergic \rpi{} cells in the anterior hindbrain that carry a
single sinusoidal activity peak travelling around the functional ring,
persisting for tens of seconds, and tracking the swum heading, supported
by the reciprocal dIPN connectivity the EM reconstruction reveals. Our
trained, connectome-locked circuit reproduces
the computational core of this picture: it integrates heading at
near-unit gain and carries a single travelling activity peak on the dIPN
ring (\cref{sec:results_evolution}), realised on the measured fish2
wiring with the magnitudes sitting on the signed connectome support. In
that sense the results agree with the introduction: the ring-attractor
account is realisable on the real connectome. They also sharpen its open
questions in two ways the introduction did not anticipate. First, the
implementation is \emph{degenerate}: many connectome-consistent circuits
integrate equally well, the preferred-heading map is not anatomically
organised, and the recurrent pool mixes heading and translation rather
than segregating them by cell type (\cref{sec:results_four_classes}),
so anatomy fixes \emph{that} the circuit integrates, not \emph{how} the
code is laid out. Second, the introductory ring activity is only
\emph{partly} matched on the recorded data: the task-only model's calcium
is nearly uncorrelated with the recording until the observation loss is
added (\cref{sec:results_obs_sweep}), and the recorded ARTR carries an
intrinsic dark-state rhythm that the command-driven model lacks
(\cref{sec:results_artr}). The agreement is therefore at the level of the
\emph{computation} the introduction describes, while the
\emph{implementation} and the \emph{full recorded dynamics} remain
under-determined, the gap the rest of the paper sets out to measure.

\begin{figure}[p]
  \centering
  \includegraphics[width=0.8\textwidth]{\figdir/fig_zebrafish_rotation_covariates.png}
  \caption{\textbf{Signals present in the recording but not given to
  the model, over the ZAPBench Rotations block} ($656$ frames
  $\approx 600\,\mathrm{s}$; shared time axis). Top to bottom: the
  periodic angular-velocity drive $s(t)$ (the model's only input),
  its integrated heading, the stimulus direction; the larval-fish
  \emph{forward} (slow, $\sim$stimulus-independent drift) and
  \emph{turning} (largely stimulus-locked with $s(t)$) swim decodes; the raw left/right tail-EMG channels (RMS)
  the decodes are computed from; the recorded $\Delta F/F$ of the three
  velocity-gate afferent classes in the \cref{fig:zhd_3d} colour code ---
  \emph{ARTR} (blue, angular drive), \emph{pt-IPN1} (orange, exteroceptive
  forward) and \emph{motor\_efferent} (green, angular efference),
  matched cells only; and two views of the recorded
  $n{=}300$ ring-pool $\Delta F/F$ --- the global ZAPBench rastermap
  permutation (matches \cref{fig:zhd_obs_baseline}) and a per-neuron
  preferred-heading-angle sort computed from the same calcium.
  These behavioural and afferent signals are the covariates discussed in
  \cref{sec:results_unmodeled} as bounding the calcium reconstruction of
  \cref{sec:results_obs_sweep}.}
  \label{fig:zhd_unmodeled}
\end{figure}

\subsection{Calcium-observation loss partially closes the activity gap}
\label{sec:results_obs_sweep}

We now turn the observation loss on. Using the sign-locked RNN of
\cref{method:RNN} (the same recurrent dynamics used throughout
Section 4) we trained the $839$-cell \emph{IPN12 precursor} circuit
(\cref{method:zapbench}; the connectome before the IPN-core families
were added to reach the nominal $917$-cell circuit) on the
swim-integration task (dataset~A) while augmenting each minibatch with
$B_c$ recorded ``Rotations'' trials (dataset~B) and matching predicted
to recorded calcium with the $z$-scored shape loss of \cref{eq:zhd_obs},
sweeping the observation weight
$\lambda_{\mathrm{obs}} \in \{0,\,0.1,\,0.5,\,1,\,5\}$
($\lambda_{\mathrm{obs}}=0$ is the task-only control). Each trained model is then driven by the
held-out recorded stimulus and its rolled-out activity rendered to
GCaMP7f calcium exactly as for the baseline (\cref{fig:zhd_obs_baseline}),
and scored against the recorded $\Delta F/F$ over a continuous
$600\,\mathrm{s}$ rollout (\cref{tab:zhd_obs_sweep},
\cref{fig:zhd_obs_reconstruction}). The figure of merit is the per-neuron
temporal $z$-MSE of \cref{eq:zhd_obs}, equal to $2(1-\rho)$ with $\rho$
the predicted-vs-recorded Pearson correlation, reported separately for
the $300$ ring-pool cells and the $181$ afferents.

The loss does what it is designed to. With $\lambda_{\mathrm{obs}} = 0$
the task-only model's calcium is essentially uncorrelated with the
recording ($z$-MSE $\approx 1.88$ for both pools, $\rho \approx 0.06$):
the network integrates heading perfectly yet its activity is unrelated to
the brain's, exactly the degeneracy of \cref{sec:results_obs_baseline}.
Turning the loss on collapses this error immediately: already at the
smallest weight $\lambda_{\mathrm{obs}} = 0.1$ the ring-pool $z$-MSE
drops to $1.25$ ($\rho \approx 0.38$) and the afferent pool to $0.92$
($\rho \approx 0.54$), and larger weights buy little further reduction
($z$-MSE plateaus at $1.1$--$1.3$ for the ring activity peak and $0.90$--$0.98$ for the
afferents across $\lambda_{\mathrm{obs}} \in [0.1, 5]$). Most of the
benefit is therefore captured by the mere presence of the term rather
than its magnitude. The afferent (\emph{RIPN}/\emph{pt-IPN}) pool is
reconstructed markedly better than the recurrent pool (afferent
$\rho \approx 0.54$ and SSIM up to $0.44$, vs.\ ring-pool $\rho \approx 0.40$
and SSIM $\approx 0.18$), consistent with the afferents carrying the
stimulus-locked part of the recording while the ring activity peak is the irregular,
drift-laden activity that is hardest to match.

The fit is partial, not complete, and trades against the task at large
weights. Even at its best the ring-pool correlation plateaus near
$\rho \approx 0.4$: the loss pulls the predicted activity toward the
recording but does not reproduce it, so the degeneracy of
\cref{sec:results_obs_baseline} is \emph{partially}, not fully, broken,
the expected outcome for a single per-neuron shape loss applied to
$481$ of the $839$ cells. Pushing $\lambda_{\mathrm{obs}}$ higher does not
help and begins to hurt: the structural SSIM peaks at the smallest weight
($\lambda_{\mathrm{obs}} = 0.1$) and declines thereafter, and the
decoded-heading panels (\cref{fig:zhd_obs_reconstruction}, bottom)
develop more frequent slips as the observation term competes with the
heading MSE. An intermediate weight
($\lambda_{\mathrm{obs}} \approx 0.1$--$0.5$) is thus the operating point
that best balances calcium realism against task accuracy.

\begin{table}[H]
\centering
\small
\begin{tabular}{c c c c c}
\toprule
$\lambda_{\mathrm{obs}}$ & ring-pool $z$-MSE & ring-pool SSIM
       & aff.\ $z$-MSE & aff.\ SSIM \\
\midrule
$0$   & $1.88$ & $\mathbf{0.29}$ & $1.83$ & $0.21$ \\
$0.1$ & $1.25$ & $0.17$ & $0.92$ & $\mathbf{0.44}$ \\
$0.5$ & $\mathbf{1.14}$ & $0.18$ & $0.98$ & $0.43$ \\
$1.0$ & $1.31$ & $0.16$ & $0.95$ & $0.37$ \\
$5.0$ & $1.21$ & $0.17$ & $\mathbf{0.90}$ & $0.22$ \\
\bottomrule
\end{tabular}
\caption{\textbf{Calcium-reconstruction quality vs.\ observation weight
$\lambda_{\mathrm{obs}}$.} GNN model, $839$-cell IPN12 circuit, scored on
the held-out continuous $600\,\mathrm{s}$ rollout against the recorded
ZAPBench Rotations block. $z$-MSE is the per-neuron temporal loss of
\cref{eq:zhd_obs} (lower is better; $z\text{-MSE} = 2(1-\rho)$ with $\rho$
the predicted-vs-recorded Pearson correlation);
SSIM is the structural-similarity index of the kinograph (higher is
better). Bold marks the best entry in each column.}
\label{tab:zhd_obs_sweep}
\end{table}

\paragraph{Task generalisation across the gcamp family.}
The observation-loss recipe is task-agnostic (the per-neuron
$z$-scored shape loss does not refer to heading, distance, or any
other behavioural variable), so the same training procedure can
be swapped onto the other swim-integration sub-tasks of
\cref{sec:results_variants}. We train two such companions of
\cref{fig:zhd_obs_reconstruction}, both supervised against the
recorded \emph{Rotations} block of dataset~B: (i) the joint
heading-plus-distance task (\cref{fig:zhd_obs_reconstruction_both},
the ``selfmotion\_both'' four-input variant) and (ii) the
translation-only task
(\cref{fig:zhd_obs_reconstruction_translation},
``selfmotion\_translation''). For the four-input model the
continuous test-time rollout is now driven jointly by the recorded
$\omega(t)$ and the recorded swim-forward envelope
$v_{\mathrm{fwd}}(t)$ (read from the calcium dataset's
swim-forward channel, the same channel the trainer feeds into the
calcium-batch trials); i.e.\ the model sees at test exactly the
four channels it was trained on. For the translation-only model
the rollout is driven by $v_{\mathrm{fwd}}$ alone. The pattern is
consistent across these two variants: the afferent pool
reconstructs at $\rho \approx 0.4$ (SSIM $0.39$--$0.48$), the
recurrent ring-pool at $\rho \approx 0.2$--$0.3$ (SSIM
$0.20$--$0.25$), and the per-frame task readouts (heading and/or
distance) remain accurately decoded in every case; the observation
loss does not collapse onto the task heads. The Turning block of
dataset~B is excluded from this comparison because its recorded
stimulus is stored as discrete L/R direction codes (no angular
velocity in the dataset's first channel), which leaves the input
gate driving the trained heading-only model with zeros and
therefore cannot supervise the heading-driven activity meaningfully.

\begin{figure}[p]
  \centering
  \includegraphics[width=0.8\textwidth]{\figdir/fig_zebrafish_calcium_reconstruction.png}
  \caption{\textbf{Calcium reconstruction with the observation loss on
  ($\lambda_{\mathrm{obs}} = 1$, heading task).} The
  sign-locked RNN of \cref{method:RNN}, trained on the
  heading-only task ($n_{\mathrm{in}} = 3$, $n_{\mathrm{out}} = 2$,
  target $(\cos\theta, \sin\theta)$) with calcium-observation
  supervision, is driven by the held-out recorded ZAPBench Rotations
  stimulus (top: angular-velocity drive $\omega(t)$) and its
  rolled-out activity rendered to GCaMP7f calcium
  (\cref{eq:zhd_gcamp}). For the recurrent pool ($300$ cells, rastermap
  row order) and each of the three afferent classes (\emph{ARTR},
  \emph{pt-IPN1}, \emph{motor\_efferent}), two rows are
  shown: the recorded $\Delta F/F$ and the model's calcium under a
  continuous $600\,\mathrm{s}$ rollout; per-panel $z$-MSE and SSIM are
  inset. Each afferent class is sorted by per-cell preferred-heading
  angle (\cref{fig:zhd_unmodeled}), the same order in its recorded and
  model rows. Bottom: true heading
  (green) and the model's readout-decoded heading (black) over the
  same continuous rollout. The task-only baseline (no observation loss)
  is \cref{fig:zhd_obs_baseline}; \cref{tab:zhd_obs_sweep} reports the
  full $\lambda_{\mathrm{obs}}$ sweep.}
  \label{fig:zhd_obs_reconstruction}
\end{figure}

\begin{figure}[p]
  \centering
  \includegraphics[width=0.8\textwidth]{\figdir/fig_zebrafish_calcium_reconstruction_both_gcamp.png}
  \caption{\textbf{Calcium reconstruction --- joint
  heading-plus-forward-distance task} ($\lambda_{\mathrm{obs}} =
  1$). Same layout as \cref{fig:zhd_obs_reconstruction} for the
  joint heading $+$ distance variant
  (``selfmotion\_both''): $n_{\mathrm{in}} = 4$, $n_{\mathrm{out}}
  = 3$, target $(\cos\theta, \sin\theta, d)$ with $d$ the
  \emph{cumulative} forward distance $d(t)=\int v_{\mathrm{fwd}}\,ds$
  (no leak --- the parent ``selfmotion\_both'' config has no
  $\xi$-decay time constant). Both drive channels the trained
  $4$-input model expects are now plotted at the top: $\omega(t)$
  is the recorded ZAPBench Rotations angular-velocity drive, and
  $v_{\mathrm{fwd}}(t)$ is the recorded swim-forward envelope read
  from the calcium dataset's $\mathrm{swim\_forward.zarr}$ (the
  same tail-EMG channel the trainer feeds into the calcium-batch
  trials, \cref{fig:zhd_unmodeled}); the continuous rollout below
  is driven by both. The bottom panel shows the model's decoded
  $d$ (black) against the cumulative reference
  $\int v_{\mathrm{fwd}}\,ds$ (green) over the continuous
  $600\,\mathrm{s}$ block.}
  \label{fig:zhd_obs_reconstruction_both}
\end{figure}

\begin{figure}[p]
  \centering
  \includegraphics[width=0.8\textwidth]{\figdir/fig_zebrafish_calcium_reconstruction_translation_gcamp.png}
  \caption{\textbf{Calcium reconstruction --- translation-only task}
  ($\lambda_{\mathrm{obs}} = 1$). Translation-only variant
  (``selfmotion\_translation''): $n_{\mathrm{in}} = 1$,
  $n_{\mathrm{out}} = 1$, target $d(t) = \int v_{\mathrm{fwd}}\,ds$
  (cumulative). The drive at the top is the recorded swim-forward
  envelope $v_{\mathrm{fwd}}(t)$ (the only channel this $1$-input
  model receives at test time); the $\omega$ panel of
  \cref{fig:zhd_obs_reconstruction_both} is absent here because the
  model has no rotation input. Below: bump-pool and afferent real /
  continuous-rollout kinographs in matched rastermap order, per-pool
  $z$-MSE / SSIM inset. Bottom: decoded forward distance $d$ (black)
  against the cumulative reference $\int v_{\mathrm{fwd}}\,ds$
  (green) over the continuous $600\,\mathrm{s}$ block; no heading
  panel since heading is not decoded in this variant.}
  \label{fig:zhd_obs_reconstruction_translation}
\end{figure}

\subsection{Proprioceptive 2-D path integration under calcium supervision}
\label{sec:results_propriocep}

The proprioception-extended circuit of \cref{method:proprioception}
(\cref{fig:zhd_3d}\textbf{d}) routes the forward-swim drive
$v_{\mathrm{fwd}}$ to IPN through two anatomically distinct ports in
parallel (the exteroceptive \emph{pt-IPN1} channel and the
motor-efferent \emph{RIPN11/12\_a/12\_c} channel) and is trained on
the full 2-D path-integration task ($n_{\mathrm{in}}=5$,
$n_{\mathrm{out}}=4$, target $(\cos\theta,\sin\theta,x,y)$) \emph{with}
the calcium-observation loss on (dataset~B ``Rotations'',
$\lambda_{\mathrm{obs}}=1$, $B_c=8$). This is the single circuit in
which the proprioceptive afferent split, the hardest (2-D path) target,
and the real-calcium supervision are all active at once.
The proprioceptive circuit differs from the task-only 2-D models
(\cref{fig:zhd_evolution_position2d,fig:zhd_evolution_position2d_leaky})
only in the second (proprioceptive) afferent port and the added
observation term. Heading is still recovered on held-out swim trials (per-trial
Pearson $r_\theta \approx 0.94$), and the 2-D-position head integrates
the path at a per-frame Euclidean RMSE of $\approx 0.38$ on natural
trials for the cumulative variant: the calcium term and the extra
afferent port do not degrade the task.

Driving these same two checkpoints with the recorded ZAPBench Rotations
stimulus and rendering the rolled-out activity to GCaMP7f calcium
(\cref{eq:zhd_gcamp}) gives the reconstructions of
\cref{fig:zhd_obs_reconstruction_propriocep_position2d,fig:zhd_obs_reconstruction_propriocep_position2d_leaky}.
The layout follows \cref{fig:zhd_obs_reconstruction_translation}, except
that the bottom panel shows the decoded 2-D \emph{trajectory}
$(\hat x,\hat y)$ against the cumulative path-integration reference, in
spatial coordinates, rather than the scalar forward-distance trace. As
elsewhere the afferent pool reconstructs markedly better than the
recurrent ring pool (cumulative: afferent $z$-MSE $=0.97$, SSIM $=0.43$
vs.\ ring-pool $z$-MSE $=1.12$, SSIM $=0.17$; leaky: afferent
$z$-MSE $=0.94$, SSIM $=0.59$ vs.\ ring-pool $z$-MSE $=1.04$,
SSIM $=0.20$), consistent with the afferents carrying the
stimulus-locked part of the recording. The decoded trajectory makes the
leaky-vs-cumulative contrast explicit on the continuous $600\,\mathrm{s}$
rollout: the cumulative decode drifts away from the origin while the
leaky decode stays bounded near it, the spatial analogue of the scalar
forward-distance trace in \cref{fig:zhd_obs_reconstruction_translation}.

\begin{figure}[p]
  \centering
  \includegraphics[width=0.8\textwidth]{\figdir/fig_zebrafish_calcium_reconstruction_propriocep_position_2d_gcamp.png}
  \caption{\textbf{Calcium reconstruction --- proprioceptive 2-D
  path-integration task, cumulative target} ($\lambda_{\mathrm{obs}} =
  1$). Same layout as \cref{fig:zhd_obs_reconstruction_translation} for
  the proprioception-extended circuit ($n_{\mathrm{in}}=5$,
  $n_{\mathrm{out}}=4$, target $(\cos\theta,\sin\theta,x,y)$). Top: the
  recorded ZAPBench Rotations angular-velocity drive $\omega(t)$ and the
  recorded swim-forward envelope $v_{\mathrm{fwd}}(t)$. For the
  ring pool ($300$ cells, rastermap order) and each of the three
  afferent classes (\emph{ARTR}, \emph{pt-IPN1}, \emph{motor\_efferent},
  sorted by per-cell preferred-heading angle), the recorded $\Delta F/F$
  and the model's GCaMP7f calcium under a continuous
  $600\,\mathrm{s}$ rollout are shown, with per-panel $z$-MSE and SSIM
  inset. \emph{Bottom panel
  (replacing the scalar-$d$ trace of
  \cref{fig:zhd_obs_reconstruction_translation}):} the decoded 2-D path
  $(\hat x,\hat y)$ (black) against the cumulative path-integration
  reference $(\,\int v_{\mathrm{fwd}}\cos\theta,\ \int
  v_{\mathrm{fwd}}\sin\theta\,)$ (green) in spatial coordinates over the
  same continuous block; the marker is the shared origin. The cumulative
  decode drifts from the origin as the unbounded reference grows.}
  \label{fig:zhd_obs_reconstruction_propriocep_position2d}
\end{figure}

\begin{figure}[p]
  \centering
  \includegraphics[width=0.8\textwidth]{\figdir/fig_zebrafish_calcium_reconstruction_propriocep_position_2d_leaky_gcamp.png}
  \caption{\textbf{Calcium reconstruction --- proprioceptive 2-D
  path-integration task, leaky target} ($\lambda_{\mathrm{obs}} = 1$,
  $\tau_p = 0.5\,$s). Same layout as
  \cref{fig:zhd_obs_reconstruction_propriocep_position2d}. In the bottom
  trajectory panel the decoded path (black) stays bounded near the
  origin --- the signature of the leaky 2-D integrator
  (\cref{eq:xy_pi}) --- while the cumulative reference (green) drifts,
  the spatial counterpart of the bounded-vs-growing distance in
  \cref{fig:zhd_obs_reconstruction_translation}.}
  \label{fig:zhd_obs_reconstruction_propriocep_position2d_leaky}
\end{figure}


\section{Cell-type vocabulary}
\label{app:celltype_vocabulary}

\Cref{tab:zhd_mapping} lists the cell-type families of the HD circuit
exposed by the \emph{fish2} neuprint server, together with their roles
and \emph{fish2} type prefixes, and is referenced throughout the main
text whenever a cell-type tag is introduced.

\begin{table}[H]
\centering
\small
\resizebox{\textwidth}{!}{%
\begin{tabular}{l l l r l}
\toprule
Role & Zebrafish population & \emph{fish2} type prefix & $n$ & Sign / source \\
\midrule
Recurrent (HD ring) & dorsal IPN (\rpi) & \emph{IPNd*} & 313 & inh.\ (Dale); \rpi{} population \\
Recurrent (HD ring, subset) & dorsal-subset IPN & \emph{IPNds*} & 95 & inh.\ (Dale) \\
Recurrent (HD ring, IPN12) & dorsal IPN & \emph{IPN12\_a/\_b} & 108 & inh.\ (v1) / exc.\ (v2) \\
Recurrent (IPN-core) & core IPN & \emph{IPN28/29/31--36} & 184 & inh.\ (Dale) \\
Angular afferent (ARTR) & Rh2/3 $\to$ IPN & \emph{RIPN01/02/03\_a/03\_b} & 76 & exc.; ARTR \citep{dunn2016artr} \\
Motor efference & motor $\to$ IPN & \emph{RIPN11/12\_a/12\_c} & 65 & exc.; efference-copy (our assignment) \\
Other RIPN (residual) & habenula $\to$ IPN & \emph{RIPN05/12\_b} & 25 & exc.; not gate-routed \\
Exteroceptive translational afferent & pretectum $\to$ IPN & \emph{pt-IPN1} & 51 & exc.; optic / water-flow (our assignment) \\
\midrule
\multicolumn{4}{l}{\emph{total}} & 917 \\
\bottomrule
\end{tabular}%
}
\caption{HD-circuit cell-type families from the \emph{fish2}
reconstruction. \emph{Role}: how the cell is wired into the model
(recurrent ring, gate-routed afferent, or residual). \emph{Population}:
the corresponding zebrafish anatomical pool. \emph{fish2 type prefix}:
the \emph{fish2} connectome type name; the asterisk is a wildcard
grouping all sub-types with that prefix (e.g.\ \emph{IPNd13B},
\emph{IPNd17A}, \emph{IPNdp01} all match \emph{IPNd*}). $n$: cell
count in this pool. \emph{Sign / source}: Dale sign used in this work
(\emph{v1} and \emph{v2} are the two IPN12 sign variants), and the
functional identification or its origin --- a citation when the pool
matches a previously named population, ``our assignment'' when the
functional label is introduced here. \emph{IPN12*} (dorsal-IPN
recurrent) and the unrelated \emph{RIPN12*} (habenular afferent) are
distinct families despite the shared number. The \emph{IPN-core}
families (\emph{IPN28/29/31--36}) complete the recurrent GABAergic
substrate; the eight rows sum to the nominal $917$-cell circuit (the
$839$-cell precursor used for the observation-loss runs of
\cref{sec:results_obs_sweep} omits IPN-core).}
\label{tab:zhd_mapping}
\end{table}

\section{Configuration families}
\label{app:config_families}

The $80$ training configurations in \texttt{config/zebrafish/} group into a
small number of families; each varies one axis (task, readout, weight prior,
connectome, stimulus, or seed) about the base heading-integration run. Most
tasks are trained both as the sign-locked recurrent network
(\texttt{zebrafish\_hd\_si\_ipn\_917\_v1\_*}) and as the message-passing GNN
(\texttt{zebrafish\_hd\_si\_gnn\_ipn\_917\_v1\_*}); the RNN/GNN columns count
the members of each. All share the nominal $917$-cell circuit unless the family
replaces the connectome. Suffixes below are relative to the
\texttt{...\_selfmotion\_rotation} base.

\begin{table}[h]
\centering\small
\begin{tabular}{@{}l c c p{7.4cm}@{}}
\toprule
Family (config suffix) & RNN & GNN & Description \\
\midrule
\multicolumn{4}{@{}l}{\emph{Task suite}}\\
\texttt{rotation}                       & 1 & 1 & Heading-only integration: $\omega\!\to\!\theta$, $\cos/\sin$ readout (base task).\\
\texttt{rotation\_vfwd}                  & 1 & 1 & Heading-only readout, but $v_{\mathrm{fwd}}$ is also fed --- latent-translation probe (\cref{app:latent_translation}).\\
\texttt{both}, \texttt{both\_leaky}      & 2 & 2 & Heading $+$ forward distance; cumulative and leaky.\\
\texttt{position\_2d}, \texttt{\_leaky}  & 2 & 0 & 2-D path integration to $(x,y)$ (the GNN 2-D runs sit under the clustered-motion family below).\\
\texttt{propriocep\_*}                   & 5 & 4 & Proprioceptive gate (sensory $\omega$ $+$ efference copy): distance, 2-D, and gain-mismatch sub-tasks.\\
\midrule
\multicolumn{4}{@{}l}{\emph{Readout ablations} (heading-only RNN)}\\
\texttt{rotation\_kbins}                 & 1 & 0 & $K{=}64$-bin heading-classifier readout (vs.\ $\cos/\sin$).\\
\texttt{rotation\_mlpdec}\,(\,$+$\,\texttt{\_\{1..5\}}) & 6 & 0 & Nonlinear MLP readout (vs.\ linear), with 5 seed replicates.\\
\midrule
\multicolumn{4}{@{}l}{\emph{Connectome weight prior} (Hulse-style block-cosine)}\\
\texttt{\{rotation,both,position\_2d\}\_cos\{000..100\}} & 15 & 0 & Anchors each (post,\,pre) cell-type block of $W^{\mathrm{rec}}$ to the connectome direction at strength $0.00$--$1.00$ (5 levels $\times$ 3 tasks).\\
\midrule
\multicolumn{4}{@{}l}{\emph{Null connectomes} (necessity controls, \cref{app:er_null})}\\
\texttt{\{rotation,rotation\_vfwd\}\_bs\_\{1..5\}} & 10 & 0 & Block-shuffle: per (post-,\,pre-type) block edge count/$|W|$/sign preserved, fine wiring randomised.\\
\texttt{rotation\_er\{,\_1..4\}}         & 1 & 4 & Erd\H{o}s--R\'enyi, size/density/$|W|$-matched.\\
\texttt{rotation\_er\_ei\_\{1..5\}}      & 0 & 5 & ER matched to the measured inhibitory fraction.\\
\midrule
\multicolumn{4}{@{}l}{\emph{Stimulus / embedding variants} (GNN)}\\
\texttt{\{rotation,position\_2d\}\_clustermove\{,\_1..5\}} & 0 & 13 & Clustered self-motion stimulus regime (bouts), with seed replicates.\\
\texttt{rotation\_freezeemb}             & 0 & 1 & Type-cluster embedding held frozen (no-learning ablation).\\
\midrule
\multicolumn{4}{@{}l}{\emph{Seed replicates} (solution variability)}\\
\texttt{rotation\_vfwd\_rep\_\{1..5\}}   & 5 & 0 & Independent random seeds of the heading\,$+\,v_{\mathrm{fwd}}$ probe.\\
\midrule
\textbf{Total} & \textbf{49} & \textbf{31} & \textbf{80 configurations}\\
\bottomrule
\end{tabular}
\caption{\textbf{Configuration families for the zebrafish HD study.}
Suffixes are relative to the \texttt{...\_selfmotion\_rotation} base (e.g.\
\texttt{both} $\equiv$
\texttt{zebrafish\_hd\_si\_ipn\_917\_v1\_selfmotion\_both}). Leaky variants
supervise a finite-$\tau$ leaky integrator instead of a cumulative one.}
\label{tab:zhd_config_families}
\end{table}

\clearpage

\bibliographystyle{unsrtnat}
\begin{thebibliography}{9}

\bibitem[Petrucco et al.(2023)]{petrucco2023}
L.~Petrucco, H.~Lavian, Y.~K.~Wu, F.~Svara, V.~\v{S}tih, R.~Portugues.
Neural dynamics and architecture of the heading direction circuit in
zebrafish. \emph{Nature Neuroscience} \textbf{26}, 765--773 (2023).

\bibitem[Allier et~al.(2026)]{cx2026}
C.~Allier et~al.
Central complex: identifiable computation, non-identifiable
implementation. Companion manuscript, in preparation (2026).

\bibitem[Dana et al.(2019)]{dana2019}
H.~Dana, Y.~Sun, B.~Mohar, et~al.
High-performance calcium sensors for imaging activity in neuronal
populations and microcompartments. \emph{Nature Methods} \textbf{16},
649--657 (2019).

\bibitem[Feierstein et al.(2023)]{feierstein2023}
C.~E.~Feierstein, M.~H.~M.~de~Goeij, A.~D.~Ostrovsky, A.~Laborde,
R.~Portugues, M.~B.~Orger, C.~K.~Machens.
Dimensionality reduction reveals separate translation and rotation
populations in the zebrafish hindbrain. \emph{Current Biology}
\textbf{33}, 3911--3925 (2023).

\bibitem[Dunn et al.(2016)]{dunn2016artr}
T.~W.~Dunn, Y.~Mu, S.~Narayan, O.~Randlett, E.~A.~Naumann, C.~T.~Yang,
A.~F.~Schier, J.~Freeman, F.~Engert, M.~B.~Ahrens.
Brain-wide mapping of neural activity controlling zebrafish exploratory
locomotion. \emph{eLife} \textbf{5}, e12741 (2016).

\bibitem[Khona \& Fiete(2022)]{khona2022}
M.~Khona, I.~R.~Fiete.
Attractor and integrator networks in the brain.
\emph{Nature Reviews Neuroscience} \textbf{23}, 744--766 (2022).

\bibitem[Seung(1996)]{seung1996}
H.~S.~Seung.
How the brain keeps the eyes still.
\emph{Proc.\ Natl.\ Acad.\ Sci.\ USA} \textbf{93}, 13339--13344 (1996).

\bibitem[Major et al.(2004)]{major2004}
G.~Major, R.~Baker, E.~Aksay, B.~Mensh, H.~S.~Seung, D.~W.~Tank.
Plasticity and tuning by visual feedback of the stability of a neural
integrator. \emph{Proc.\ Natl.\ Acad.\ Sci.\ USA} \textbf{101},
7739--7744 (2004).

\bibitem[Gon\c{c}alves et al.(2014)]{goncalves2014}
P.~J.~Gon\c{c}alves, A.~B.~Arrenberg, B.~Hablitzel, H.~Baier,
C.~K.~Machens.
Optogenetic perturbations reveal the dynamics of an oculomotor
integrator. \emph{Frontiers in Neural Circuits} \textbf{8}, 10 (2014).

\bibitem[Kim et al.(2017)]{kim2017}
S.~S.~Kim, H.~Rouault, S.~Druckmann, V.~Jayaraman.
Ring attractor dynamics in the \emph{Drosophila} central brain.
\emph{Science} \textbf{356}, 849--853 (2017).

\bibitem[Chaudhuri et al.(2019)]{chaudhuri2019}
R.~Chaudhuri, B.~Ger\c{c}ek, B.~Pandey, A.~Peyrache, I.~Fiete.
The intrinsic attractor manifold and population dynamics of a canonical
cognitive circuit across waking and sleep. \emph{Nature Neuroscience}
\textbf{22}, 1512--1520 (2019).

\bibitem[Lyu et al.(2022)]{lyu2022}
C.~Lyu, L.~F.~Abbott, G.~Maimon.
Building an allocentric travelling direction signal via vector
computation. \emph{Nature} \textbf{601}, 92--97 (2022).

\bibitem[Yang et al.(2024)]{yang2024}
C.~Yang, L.~Mammen, B.~Kim, M.~Li, D.~N.~Robson, J.~M.~Li.
A population code for spatial representation in the zebrafish
telencephalon. \emph{Nature} \textbf{634}, 397--406 (2024).

\end{thebibliography}

\end{document}
